% ============================================================
% Elsevier TUST manuscript template
% Journal: Tunnelling and Underground Space Technology
% ============================================================
\documentclass[preprint,12pt,authoryear]{elsarticle}

% --- Packages ---------------------------------------------------------
\usepackage[T1]{fontenc}
\usepackage[utf8]{inputenc}
\usepackage{amsmath,amssymb,amsthm}
\usepackage{graphicx}
\usepackage{booktabs}
\usepackage{tabularx}
\usepackage{longtable}
\usepackage{array}
\usepackage{multirow}
\usepackage{calc}
\usepackage{textcomp}
\usepackage{textcomp}
\usepackage{svg}
\usepackage{caption}
\usepackage{subcaption}
\usepackage{enumitem}
\usepackage{setspace}
\usepackage{url}
\usepackage{xurl}
\usepackage{algorithm}
\usepackage{algpseudocode}
\usepackage{listings}
\usepackage{xcolor}
\usepackage{threeparttable}
\usepackage{amssymb}
\usepackage{lineno}
\usepackage{hyperref}
\usepackage{geometry}
\usepackage{float}
\usepackage{subcaption}
\geometry{a4paper, top=2cm, bottom=2cm, left=2cm, right=2cm}
\onehalfspacing
\linenumbers
\hypersetup{colorlinks=false}
\usepackage{titlesec}
\titleformat{\subsection}{\normalfont\bfseries}{\thesubsection}{1em}{}
% Tell the svg package where the figures live
\svgpath{{Fig_SVG/}}
\graphicspath{{Fig_PDF/}}
\AtBeginEnvironment{table}{\small}
\AtBeginEnvironment{longtable}{\small}

% --- Pandoc compatibility --------------------------------------------
\providecommand{\tightlist}{%
  \setlength{\itemsep}{0pt}\setlength{\parskip}{0pt}}
\setlength{\emergencystretch}{3em}
\makeatletter
\def\maxwidth{\ifdim\Gin@nat@width>\linewidth\linewidth\else\Gin@nat@width\fi}
\def\maxheight{\ifdim\Gin@nat@height>\textheight\textheight\else\Gin@nat@height\fi}
\makeatother
\setkeys{Gin}{width=\maxwidth,height=\maxheight,keepaspectratio}
\journal{Tunnelling and Underground Space Technology}

% --- Theorem-like environments ---------------------------------------
\newtheorem{definition}{Definition}
\newtheorem{property}{Property}

% --- Listing style for SPARQL / pseudocode ---------------------------
\lstdefinelanguage{SPARQL}{
  morekeywords={SELECT,WHERE,PREFIX,FILTER,OPTIONAL,UNION,DISTINCT,ORDER,BY,LIMIT,a},
  sensitive=true,morecomment=[l]{\#},morestring=[b]"
}
\lstset{language=SPARQL,basicstyle=\ttfamily\small,
        frame=single,breaklines=true,xleftmargin=8pt}

\begin{document}

\begin{frontmatter}

\title{A Prescriptive Digital-Twin Framework for Shield Metro Tunnels: Multimodal Point-Cloud Perception and Standards-Grounded Reasoning for Maintenance Decision Support}

\author[monash]{Mohammad Matin Rouhani}
\author[swinburne]{Mengqi Huang}
\author[monash]{Qian-Bing Zhang\corref{cor1}}
\ead{qianbing.zhang@monash.edu}
\cortext[cor1]{Corresponding author}

\affiliation[monash]{organization={Department of Civil and Environmental Engineering, Monash University}, city={Clayton}, country={Australia}}
\affiliation[swinburne]{organization={Department of Civil and Construction Engineering, Swinburne University of Technology}, city={Melbourne}, country={Australia}}

\begin{abstract}
Underground transport networks carry decades of inspection and rehabilitation liabilities, yet 
most tunnel digital twins fall short of prescribing maintenance, focusing instead on describing 
or forecasting conditions. This paper proposes a prescriptive multimodal digital-twin pipeline 
for segmental shield-tunnel maintenance, which consists of three stages: point-cloud leakage 
segmentation, geometrical deformation evaluation, and an ontology-based decision layer. Four 
lining classes (leakage, joints, segments, and pockets) are segmented using Sonata and evaluated 
under two transfer regimes and three label fractions with and without a supplementary Australian 
reference tunnel. When fully fine-tuned on the in-domain Nanjing Metro data with full 
supervision (10\% of the scanned dataset), it reached a leakage Intersection-over-Union of 94.34 
and a mean accuracy of 98.60, outperforming nine reported point-cloud baselines, while a leakage 
IoU of 87.59 was retained with only a quarter of the labels; cross-site data was most useful in 
the low-label, frozen-encoder setting. A multi-zone polynomial fit was used to reconstruct each 
ring, and these were reduced to offset-immune deformation indicators. Ovalization was cross-
validated against an independent conic fit, yielding a mean absolute difference of 1.37 mm. 
Using schema-guided large-language-model extraction, we populated 51 failure-mode reasoning 
chains from the governing Chinese urban-rail standards, with every ring graded and provenance 
retained to the source clause. Applied to two held-out tunnels, the pipeline generated 
traceable, standard-referenced maintenance prescriptions, providing a label-efficient route from 
raw scan to defensible decision.
\end{abstract}

\begin{keyword}
Point cloud \sep Semantic segmentation \sep Digital twin \sep Tunnel maintenance \sep Point Transformer v3 \sep Leakage detection
\end{keyword}

\end{frontmatter}

\section*{Highlights}
\begin{itemize}
\item 94.3\% leakage IoU, the top value across all compared point-cloud models in previous studies.
\item 50\% of labels (approximately 5\% of scans) beat the prior best leakage IoU.
\item Competitive 87.6\% leakage IoU from just 25\% of labels (approximately 2.5\% of scans).
\item 51 standard-derived FMEA rules yield clause-traceable ring decisions.
\end{itemize}

\section{Introduction}\label{introduction}

Underground transport infrastructure is expanding rapidly, and every new
tunnel commissioned today carries inspection, monitoring, and
rehabilitation liabilities spanning decades of service life. The real
question for asset owners is not whether to maintain, but how to
allocate finite maintenance effort across a growing portfolio of ageing
assets with timely and traceable condition evidence. Building
Information Modelling (BIM) partly backstops this need through as-built
documentation, spatial indexing of inspection records, and work-order
management. But it doesn't, by itself, translate condition observations
into maintenance decisions. The digital twin (DT) extends BIM with live
data connectivity and, more importantly for maintenance, a reasoning
layer that connects the sensed state of the structure to engineering
knowledge ~\citep{Grieves2017, Tao2019a}. Digital twin
(DT) capabilities are typically categorised into four levels based on a
maturity model: descriptive twins describe the current state, reflective
(diagnostic) twins analyse historical trends and design intent,
predictive twins forecast the future deterioration, and prescriptive
twins determine appropriate actions based on combining condition
assessments and intervention knowledge ~\citep{Babanagar2025,Zhu2025}. 
Looking into recent tunnel-DT studies, a noticeable maturity
gradient emerges: the field predominantly focuses on descriptive and
predictive levels, while truly prescriptive twins, those that provide a
traceable maintenance prescription rather than merely a forecast, are
relatively scarce and largely remain conceptual.

This imbalance reflects where the demanding step lies. Reaching the
prescriptive level requires that heterogeneous condition observations be
connected to a structured body of maintenance knowledge through which
decision logic can be traversed automatically, and the reviewed
literature shows that most frameworks stop short of this. Many propose
architectures for predictive maintenance ~\citep{Khan2025, Zhao2024}
or demonstrate high-fidelity reconstruction and visualisation ~\citep{Cui2026, Fang2024}, and a growing number target operational
optimisation for tunnel-boring machines ~\citep{Yang2026}. Relatively
few close the loop from sensed defect to prescribed intervention for an
operating lining. The most frequently reported limitation across the
reviewed corpus is data integration and interoperability, followed by
the unavailability of high-quality labelled data, and in particular,
failure data, as most available records describe structures in normal
condition. A recurring corollary is that complex defects, such as cracks
and leakages, are difficult to model geometrically and that the matching
of each data type to an appropriate algorithm is still unresolved. The
collective picture is one of a field rich in predictive frameworks but
segmented from practical, fully integrated prescriptive implementation.

These obstacles motivate a multimodal approach to condition assessment.
No single inspection modality captures the full state of a deteriorating
lining: surface imagery and laser-scanned point clouds reveal the
visible and geometric manifestation of a defect, while subsurface and
internal methods reveal its extension and driver, so that different
modalities probe different stages of the same physical process ~\citep{He2024, Zhu2025}. Yet classical signal-level fusion of these
streams produces a sharper composite picture without producing a
maintenance prescription: combining the data is necessary but not
sufficient. The process of generating prescription necessitates a
comprehensive interpretation of the outputs from various modalities,
underpinned by engineering knowledge that links observed phenomena to
appropriate actions. This interpretation requires the application of an
ontology-based reasoning framework, which is consistently identified as
a critical yet absent component in the reviewed literature. The present
work is positioned at this gap. It develops a prescriptive, multimodal
digital-twin pipeline for shield-tunnel maintenance in which point-cloud
and image evidence are interpreted through a structured,
standards-referenced knowledge base, and it is organised around the
three stages that this pipeline integrates: leakage prediction,
deformation evaluation, and an ontology-based decision layer.

\section{Related works}
\label{related-works}

\subsection{\textbf{Leakage Prediction}}%\label{leakage-prediction}

Automated leakage assessments in shield tunnels have shifted decisively
from image-based inspection toward direct interrogation of
three-dimensional scan data, and the reviewed studies divide cleanly
along that line. Early and still-common approaches segment leakage on
two-dimensional imagery: fully convolutional networks applied to
mobile-LiDAR intensity images ~\citep{Cheng2021}, attention-augmented
U-Net variants ~\citep{Wang2024_2}, KAN-based lightweight segmentation
networks ~\citep{He2025}, and instance-segmentation models built on
Mask R-CNN and BlendMask backbones ~\citep{Geng2023, Wang2024}. The tunnel environment itself, like low and unstable
illumination, is the biggest robustness in 2D image studies. This issue
has motivated explainable infrared-RGB fusion networks that recover
leakage targets invisible in ordinary photographs ~\citep{Yao2025} and
dedicated detectors that synthesize low-light conditions during training
~\citep{Liu2024}.

A second, faster-growing branch operates directly on point clouds, which
preserves the metric geometry needed to locate and size a leak on the
lining surface. Within this branch the methods span classical and
rule-based pipelines through transformer architectures. Geometric and
rule-based extraction quantifies leakage location and area from raw
laser-scan points ~\citep{Li2024}, while unsupervised schemes such
as 3D Otsu thresholding combined with adaptive denoising and KNN
clustering detect seepage without labelled data ~\citep{Wang2024}.
Learning-based segmentation has progressed from classical machine
learning on engineered descriptors ~\citep{Huang2026a} to deep
architectures purpose-built for tunnel leakage, including the ensemble
sparse-U-Net/transformer model BBE-Net ~\citep{Jiang2025b} and the
geometry-enhanced point transformer GEPT-Net ~\citep{Jiang2025a}.
Hybrid strategies bridge the two modalities by projecting points to
images: panoramic-image Mask R-CNN with deformable attention ~\citep{Li2026} and YOLO-based segmentation of LiDAR-derived greyscale images
~\citep{Liang2025} both exploit mature 2D detectors while retaining a
path back to 3D coordinates. The summary of most of the studies on
leakage predictions are presented in Table 1.

The present work sits within the point-cloud branch. It performs
multi-class segmentation on the native 3D scan rather than on projected
imagery, and it treats detected leakage not as a terminal output but as
graded evidence feeding a maintenance-decision stage, addressing the
type-distinction and risk-quantification limitations that the reviewed
detection-only methods leave open.

%\emph{Table 1- A summary of previous studies}

\begin{longtable}[]{@{}
  >{\raggedright\arraybackslash}p{(\columnwidth - 8\tabcolsep) * \real{0.1728}}
  >{\raggedright\arraybackslash}p{(\columnwidth - 8\tabcolsep) * \real{0.2036}}
  >{\raggedright\arraybackslash}p{(\columnwidth - 8\tabcolsep) * \real{0.2200}}
  >{\raggedright\arraybackslash}p{(\columnwidth - 8\tabcolsep) * \real{0.2357}}
  >{\raggedright\arraybackslash}p{(\columnwidth - 8\tabcolsep) * \real{0.1678}}@{}}
\caption{- Structural-health degree and the corresponding maintenance
measure ~\citep{MOHURD2018}}\label{Tab:01}
\tabularnewline
\toprule\noalign{}
\begin{minipage}[b]{\linewidth}\raggedright
\textbf{Study}
\end{minipage} & \begin{minipage}[b]{\linewidth}\raggedright
\textbf{Input modality}
\end{minipage} & \begin{minipage}[b]{\linewidth}\raggedright
\textbf{Method (model)}
\end{minipage} & \begin{minipage}[b]{\linewidth}\raggedright
\textbf{Method class}
\end{minipage} & \begin{minipage}[b]{\linewidth}\raggedright
\textbf{Task}
\end{minipage} \\
\midrule\noalign{}
\endfirsthead
\toprule\noalign{}
\begin{minipage}[b]{\linewidth}\raggedright
\textbf{Study}
\end{minipage} & \begin{minipage}[b]{\linewidth}\raggedright
\textbf{Input modality}
\end{minipage} & \begin{minipage}[b]{\linewidth}\raggedright
\textbf{Method (model)}
\end{minipage} & \begin{minipage}[b]{\linewidth}\raggedright
\textbf{Method class}
\end{minipage} & \begin{minipage}[b]{\linewidth}\raggedright
\textbf{Task}
\end{minipage} \\
\midrule\noalign{}
\endhead
\bottomrule\noalign{}
\endlastfoot
~\citep{Cheng2021} & LiDAR intensity images & FCN & CNN
encoder--decoder & Segmentation \\
~\citep{Geng2023} & Lining images & Improved BlendMask & Instance seg.
(DL) & Segmentation \\
~\citep{Li2024} & 3D laser-scan points & Geometric/rule-based
software & Rule-based & Measurement/ quantification \\
~\citep{Wang2024_2} & Images & CSG-UNet (attention) & CNN
encoder--decoder & Segmentation \\
~\citep{Liu2024} & Images & Tunnel Defects Detector (TDD) & Object
detection (DL) & Detection \\
~\citep{Wang2024} & 3D point clouds & 3D Otsu-KNN (unsupervised) &
Classical ML & Detection \\
~\citep{Wang2024_3} & Tunnel leakage images & Cascade-MRegNetX & Mask
R-CNN & Segmentation \\
~\citep{Wang2025} & Images & Semi-supervised teacher--student &
Semi-supervised & Segmentation \\
~\citep{Jiang2025b} & 3D point clouds & BBE-Net (sparse U-Net +
transformer) & Deep learning (ensemble) & Segmentation \\
~\citep{Jiang2025a} & Point clouds & GEPT-Net & Point transformer &
Segmentation \\
~\citep{Liang2025} & LiDAR → greyscale images & YOLOv11m-seg & YOLO &
Segmentation \\
~\citep{He2025} & Lining images & EMS-UKAN & KAN-based DL &
Segmentation \\
~\citep{Yao2025} & Infrared + RGB images & EG-UNet (explainable) & CNN
encoder--decoder & Segmentation \\
~\citep{Zhang2025} & 3D points + 2D projection & SAM + U-Net + DGCNN &
Foundation model & Segmentation \\
~\citep{Huang2026b} & MLS point clouds & Intensity-aware features +
CatBoost & Classical ML & Segmentation \\
~\citep{Li2026} & Laser scan → panoramic images & BMask R-CNN-DAT &
Mask R-CNN & Segmentation \\
\end{longtable}

\subsection{\textbf{Deformation Evaluation}}\label{deformation-evaluation}

Point cloud-based deformation evaluation of shield tunnels has evolved
along a clear path, from purely geometric fitting toward learning-based
segmentation that first isolates the structural elements and only then
measures their distortion. The geometric branch remains the workhorse
for cross-sectional metrics: The Iterative Closest Point algorithm was
combined with ellipse fitting on mobile-laser-scanning data to detect
cross-section displacement and shape change ~\citep{Camara2025}, and
the deformed ring was reconstructed through polynomial axis extraction
and binary-image processing to recover both section deformation and
segment dislocation ~\citep{Zhang2024}. These methods are accurate and
label-free, but the same fragility was repeatedly identified: the
fitting assumptions degrade in noisy, occluded, or irregular scenes.

The second branch, in which the lining is segmented into rings,
segments, and joints by a network so that deformation is computed on
clean, semantically isolated geometry. The trend is visible across a
tightly connected group of recent studies: GL-Net, a global-feature
segmentation network for shield-tunnel structural elements, was
introduced ~\citep{Li2023}. Segmentation networks (TPSegNet, then the
coordinate-only MTPointNet) were coupled with cascaded geometric schemes
to detect longitudinal and convergence deformation from robot- and
quadruped-acquired scans ~\citep{Feng2025,Feng2026}. Joint-level
deformation is treated by the same logic: two CNNs were combined with
RANSAC plane fitting to quantify longitudinal dislocation ~\citep{Yu2021}, 
ring and longitudinal joints were identified through Hough
transform and MSAC on a cylindrical projection ~\citep{Shen2025}, and
points were converted to depth and grayscale images for multi-modal
YOLOv8n quantification of joint opening ~\citep{Liu2025}. Across
these, segmentation is consistently the enabling step rather than the
end goal, since the measured indicator (convergence, dislocation,
opening) is derived geometrically once the elements are separated.

This study uses the segmentation-then-measure paradigm that the
literature has converged on, but advances it on two fronts that the
reviewed work leaves open. It carries the segmented elements through to
a quantitative, standard-referenced deformation grade, addressing the
measurement-and-assessment gap left unresolved by detection-oriented
studies.

\subsection{\textbf{Ontology-Based Decision}}\label{ontology-based-decision}

Reaching the prescriptive level of a maintenance digital twin requires a
formal knowledge layer through which condition observations can be
linked to intervention decisions and traversed automatically, and
ontologies provide that layer. Distributed ontology-based
interoperability, in which purpose-built domain ontologies are linked
through Web Ontology Language (OWL) alignment axioms, offers modular
flexibility relative to fixed-schema approaches, and has been
demonstrated across civil infrastructure, including ontological
knowledge bases for bridge maintenance ~\citep{Ren2019}, building
monitoring ~\citep{Mahdavi2017}, and integrated
geospatial-model-knowledge representation for railway digital twins ~\citep{Li2022}. The reasoning capability such structures afford depends
on the semantic richness of their model layer ~\citep{Tao2019b}, which
is precisely what distinguishes a knowledge-bearing decision layer from
a passive data repository.

In the tunnel-maintenance domain specifically, semantic-web technologies
have been applied to fuse condition information across data, object, and
knowledge levels: the COBie standard was extended ontologically to
organise tunnel digital-twin data, and a rule-based semantic-reasoning
approach was introduced for equipment fault-cause analysis ~\citep{Yu2021}, and a knowledge-driven decision-support method for tunnel
electromechanical-equipment maintenance was subsequently developed by
the same group ~\citep{Yu2023}. An ontology- and graph-based modelling
paradigm built from five elements, including: scenario, virtual model,
physical entity, relation, and component, was proposed to support
multidisciplinary operation across foundation excavation, subway
tunnels, and pipe networks ~\citep{Li2024_2}, while lifecycle
underground-utility information was modelled ontologically to support
operation and maintenance ~\citep{Wang2021}. A complementary thread concerns
populating such ontologies from the standards themselves: large language
models perform well on structured information extraction from technical
and regulatory documents, and schema-guided extraction, which constrains
the model to a predefined output structure, achieves greater
terminological consistency than free-form extraction and yields output
directly compatible with ontology population ~\citep{Caufield2024, Dagdelen2024}.
In a review of this trajectory, the maintenance-knowledge integration
layer has been identified as the principal barrier to the
predictive-to-prescriptive transition for underground infrastructure
~\citep{Zhu2025}.

This study develops an integration layer of shield-tunnel maintenance
digital twins. A schema-guided large-language-model extraction pipeline
is employed to build a domain knowledge base from the governing Chinese
urban-rail maintenance standards, representing each rule as a
failure-mode-and-effects-analysis reasoning chain that connects a
degradation mechanism, through a measurable indicator and a graded
threshold, to a prescribed intervention. The measured evidence of the
current pipeline is then scored against this extracted rule base so that
each ring-level maintenance decision is derived by traversing the
knowledge chain and is traceable to the specific standard clause that
governs it. In this way, the decision stage connects the multimodal
condition assessment with the ontology-based reasoning presented here,
closing the prescriptive loop from the sensed defect to the
standard-referenced maintenance action.

\section{Database}\label{database}

\subsection{\textbf{Australian Reference Tunnel
(ART)}}\label{australian-reference-tunnel-art}

The Australian Reference Tunnel was a small fraction (13 rings) of the
typical metro tunnels in Australia, specifically in Melbourne, Victoria,
with a 6.3m clearance diameter and 6 segment-per-ring. The scanning
process was done using Trimble X7 using 3 scans. Since the
number of scans per ring was much higher than in the other studies, the
given dataset is expected to improve the prediction stage by providing
more accurate details for each component. Also, there has been zero
leakage reported in this tunnel; thus, this tunnel can only support the
study with the other three components: Joints, Pockets, and Segments.

%\begin{figure}[htbp]
%  \centering
%  \includegraphics[width=0.92\textwidth]{Fig_PDF/Picture01.pdf}
%  \caption{ Point cloud model of the TBM tunnel: (a) scan view, (b) global view, and (c) segmental view ~\citep{ZHUn.d.}}
%  \label{fig:Fig1}
%\end{figure}




\subsection{\textbf{Nanjing metro Line 2
(S3DIS)}}\label{nanjing-metro-line-2-s3dis}

S3DIS\_Leakage dataset (S3Ld) is a publicly reported benchmark for
water-leakage segmentation in operational shield tunnels ~\citep{Jiang2025a, Jiang2025b}. The data were acquired with a FARO Focus 330
laser-scanning system along the operational shield tunnels between
Jiqingmen Street and Xinglong Street stations on Line 2 of the Nanjing
Metro, a section of approximately 1.6 km situated in the alluvial plain
of the lower Yangtze River. From almost 90 tunnel sets in the
given dataset, 9 of them have been annotated and used in this study to
evaluate the size effect during prediction.


%\begin{figure}[htbp]
%  \centering
%  \includegraphics[width=0.92\textwidth]{Fig_PDF/Picture02.pdf}
%  \caption{ Route of Nanjing Metro Line 2 ~\citep{Gong2023}}
%  \label{fig:Fig2}
%\end{figure}

\subsection{Data preprocessing}\label{data-preprocessing}

Before the point clouds were partitioned for training, the raw scans
were imported into CloudCompare for cleaning. Two operations were
carried out. First, scanning noise and spurious returns, isolated points
and low-density outliers produced by reflective lining surfaces,
ancillary equipment, and acquisition artefacts, were removed, so that
the subsequent geometric fitting and semantic segmentation operated on
the lining surface rather than on measurement noise. Second, the
residual points belonging to the tunnel invert, the track-bed region at
the base of the bore, were segmented out, leaving the exposed
segmental-ring surface over an open angular arc; this is the open-invert
geometry that the per-ring deformation reconstruction subsequently
accounts for. The cleaned point clouds, which carry the point-level
class labels (leakage, joints, segments, and pockets) used downstream,
were then exported for the data-separation stage.

\section{Methodology}\label{methodology}

The pipeline introduced is named TunSPEC  (Tunnel Standards-grounded Prescriptive 
Evaluation of Conditions). The acronym tracks its structure: two parallel point-
cloud branches carry out the Evaluation of Conditions, and a Standards-grounded l
ayer renders a Prescriptive maintenance decision from them. 
TunSPEC begins with data collection from two different sites. A
Chinese metro tunnel and an Australian Reference Tunnel (Explained
earlier in section \ref{database}), to evaluate that what fracture of the dataset can
reach an acceptable performance through a foundation model, the given
S3DIS dataset (from Nanjing Metro Line 2) was considered with three
different levels: 25\% of the whole annotated dataset, 50\%, and 100\%.
It is worth mentioning again that the whole annotated dataset was only
10\% of the whole scanned dataset; thus, the full case scenario had only
seen 10\% of previous studies.

The Evaluation-of-Conditions stage holds the two parallel condition-assessment
branches  (Figure \ref{fig:Fig3}). The leakage branch passes the point cloud through a
self-supervised semantic-segmentation model, Sonata, assigning each
point to one of the four classes. The deformation branch operates ring
by ring: each ring is aligned to its principal axis by
principal-component analysis (PCA), fitted with a robust ellipse or
circle and a multi-zone polynomial, and reduced to a set of per-ring
geometry indicators. The two branches share the same conditioned point
cloud but produce evidence of two different kinds, one perceptual and
one geometric, that the decision level later reads together.

Standards-grounded, prescriptive stage comprises the AI-assistant decision layer of 
Stage 5. In that step, the maintenance standards collected in Stage 1 are passed
to a large language model (Claude Opus 4.8 ~\citep{Anthropic2026}) under a schema-
guided extraction that returns a structured FMEA rule base. At inference, the
decision layer reads the leakage and deformation measurements of Stage
4, matches each measurement against the extracted standard thresholds,
assigns a five-level health grade, and maps that grade to a maintenance
decision with provenance retained to the source standard clause that
governed it.

\begin{figure}[htbp]
  \centering
  \includegraphics[width=0.92\textwidth]{Fig_PDF/Picture033.pdf}
  \caption{ The overall framework of this study}
  \label{fig:Fig3}
\end{figure}


\subsection{\textbf{Sonata}}\label{sonata}

The tunnel lining point clouds were segmented into four classes
(Leakage, Joints, Segments, and Pockets) using Sonata, a self-supervised
point-cloud representation model released by Meta AI ~\citep{Wu2025}.
The overall network is shown in  (Figure \ref{fig:Fig4}) (a). It comprises three parts: a pre-trained hierarchical encoder, a non-parametric upcasting stage to
restore the per-point resolution, and a lightweight segmentation head.
The encoder is the serialisation-based point transformer on which Sonata
is pre-trained ~\citep{Wu2024}. Sonata converts the unstructured
points into ordered sequences along space-filling curves and processes
them through a point-embedding stem followed by four hierarchical stages
(Figure \ref{fig:Fig4} (a)). Each stage downsamples the cloud and doubles the feature
width, taking the representation from full resolution to 1/2, 1/4, 1/8,
and 1/16 with channel widths 2C, 4C, 8C, and 16C, the last forming the
bottleneck. Each stage is built from serialised-attention blocks whose
internal structure is detailed in Figure \ref{fig:Fig4} (b): an enhanced conditional
positional encoding (xCPE), utilized as a depth-wise sparse convolution
applied through a residual connection, precedes a pre-normalised
serialised self-attention operating over shifted point patches, which is
followed by a pre-normalised feed-forward MLP (linear→GELU→linear), each
sub-layer wrapped in a residual connection ~\citep{Wu2024}.
Downsampling between stages uses grid pooling (Figure \ref{fig:Fig4} (c)): points are
assigned to voxels at the pooling grid size, their features are
max-pooled within each occupied cell, and a linear projection doubles
the channel dimension (C→2C) to yield the coarser point set ~\citep{Wu2024}. 
Each point is presented to the encoder as a nine-dimensional
feature of metric coordinates, a colour value, and a k-nearest-neighbour
surface normal (k = 15) ~\citep{Wu2025}. Sonata is an encoder-only
model, so dense per-point features are recovered by the training-free
upcasting of Figure 1 (a, d) rather than by a learned decoder ~\citep{Wu2025}. The only learned parameters beyond the bottleneck are in the
segmentation head (Figure \ref{fig:Fig4} (a)), a single linear layer followed by a
SoftMax that maps each upcast point feature to the four class logits;
the label of every original point is then recovered through the stored
inverse indices, so that all metrics are computed at full scan
resolution.

\begin{figure}[htbp]
  \centering
  \includegraphics[width=0.92\textwidth]{Fig_PDF/Picture04.pdf}
  \caption{ the architecture of Sonata (a) overview, (b)
serialised-attention block, (c) grid pooling, (d) upcasting}
  \label{fig:Fig4}
\end{figure}

The network was trained under a single fixed recipe to enable controlled
comparison. For the optimisation stage, AdamW with focal loss ($\gamma = 2$)
was used. Geometric augmentation was applied only during training. Two
transfer-learning regimes were compared: full fine-tuning, in which the
entire pre-trained encoder is updated with a reduced backbone learning
rate alongside the head, and linear probing, in which the encoder is
frozen and only the linear head is trained --- the latter directly
assessing the quality of the self-supervised representation. To
characterise the dependence on annotation quantity, each regime was run
at three labelled-data fractions (100\%, 50\%, and 25\% of the training
source files, drawn by a fixed seed and stratified by class so that
every configuration trains on a nested subset of the same data); the
validation and test sets were never reduced. The model was trained on
the labelled rings of seven tunnels and evaluated on the two fully
held-out tunnels, with the held-out split enforced at the tunnel level
to prevent leakage between training and test. The contribution of the
supplementary Australian Reference Tunnel (ART) structural dataset was
isolated by repeating every configuration with and without ART in the
training pool. In contrast, the ART data were confined to training and
validation and never entered the test set. Performance was reported per
class and overall using accuracy, precision, recall, F1,
Intersection-over-Union, Cohen's $\kappa$, and the Matthews correlation
coefficient, all computed on the original full-resolution test points.

\subsection{Multi-Zone Polynomial}\label{multi-zone-polynomial}

The reconstruction follows the multi-zone polynomial scheme of ~\citep{Lin2025}, 
but with some changes. Whereas Lin et al. normalize each
ring by fitting an elliptic cylinder to recover the central axis (their
Algorithm 1) and localize the longitudinal joints from the known, fixed
segment-installation geometry by optimizing the key-segment reference
angle (their Algorithm 2), the present method estimates the axis by
principal-component analysis (PCA), determines the cross-sectional
centre by a robust constrained-radius circle fit, and localizes joints
empirically from the independently labelled joint point set using an
anisotropy-weighted angular histogram, with an additional
largest-empty-arc detection and angular re-referencing introduced to
handle the removed invert

The evaluation is per ring and for each ring, 5 stages have been
applied, as follows: (i) estimation of the local tunnel axis and rigid
alignment of the ring point cloud to a unified cross-sectional
coordinate system (UCS); (ii) robust determination of the ring center by
a constrained-radius circle fit; (iii) reduction to polar coordinates
with detection and handling of the removed invert; (iv) localization of
longitudinal joints and recovery of the individual segment instances;
and (v) reconstruction of the radial profile of each segment by a
multi-zone polynomial fit, from which the deformation indicators (radial
deviation, joint dislocation, joint rotation, diametric convergence and
opalization)

\subsubsection{Tunnel-axis estimation and alignment to the
UCS}\label{tunnel-axis-estimation-and-alignment-to-the-ucs}
\addcontentsline{toc}{paragraph}{Tunnel-axis estimation and alignment to
the UCS}

the segment points of a ring are defined as
\(\left\{ P_{i} \right\}_{i = 1}^{N}\)in which
\(P_{i} = \left( x_{i},\ y_{i},\ z_{i} \right)^{T}\). The axis is
defined as the direction of least spatial variance, rooted from the
principal-component decomposition of a reference ring:


\begin{equation}
  \overline{P} = \frac{1}{N}\sum_{i=1}^{N}P_{i}
  \label{eq:01}
\end{equation}

\begin{equation}
  C = \frac{1}{N}\sum_{i=1}^{N}\left( P_{i} - \overline{P} \right)\left( P_{i} - \overline{P} \right)^{T}
  \label{eq:02}
\end{equation}

With the eigenpairs \(Ce_{k} = \lambda_{k}e_{k}\)ordered
\(\lambda_{1} \leq \lambda_{2} \leq \lambda_{3}\), The unit axis is the
eigenvector of the smallest eigenvalue, sign-oriented to the positive
global x direction:


\begin{equation}
  \widehat{a} = sgn\left( e_{1,x} \right)e_{1},
  \label{eq:03}
\end{equation}

The ring is then rotated so that \(\widehat{a}\) coincides with the
global axis \(\widehat{X} = (1,0,0)^{T}\). Using Rodrigues' construction
with \(V = \widehat{a} \times \widehat{X}\) so:
\(s = \left\| V \right\|\ \&\ c = \widehat{a}.\widehat{X}\),


\begin{equation}
  R = I + \lbrack V\rbrack_{\times} + \lbrack V\rbrack_{\times}^{2}\frac{1 - c}{s^{2}}
  \label{eq:04}
\end{equation}
\begin{equation}
  \lbrack V\rbrack_{\times} = \begin{bmatrix}
0 & - v_{3} & v_{2} \\
v_{3} & 0 & - v_{1} \\
 - v_{2} & v_{1} & 0
\end{bmatrix}\
  \label{eq:05}
\end{equation}



Each point is mapped to the UCS by:

\begin{equation}
  {P'}_{i} = R\left( P_{i} - \overline{P} \right) = \left( {x'}_{i},\ {y'}_{i},\ {z'}_{i} \right)^{T}\
  \label{eq:06}
\end{equation}

So, x' goes with the tunnel axis and (y', z') spans the cross-sectional
plane.

\subsubsection{Ring-centre determination by robust constrained-radius
circle
fitting}\label{ring-centre-determination-by-robust-constrained-radius-circle-fitting}
\addcontentsline{toc}{paragraph}{Ring-centre determination by robust
constrained-radius circle fitting}

Since the invert section of the tunnel was flat and then removed from
the point cloud, an unconstrained circle/ellipse fit to the partial arc
is poorly defined. Thus, the centre of the cross-section
\(\left( y_{c},\ z_{c} \right)\) is derived from a fixed design radius
to reduce the robust (soft-\(L_{1}\)) loss over the geometric residuals


\begin{equation}
  [f_{i}\left( y_{c},\ z_{c} \right) = \sqrt{\left( \ {y'}_{i} - y_{c} \right)^{2} + \left( \ {z'}_{i} - z_{c} \right)^{2}} - R\
  \label{eq:07}
\end{equation}


\begin{equation}
  \left( {\widehat{y}}_{c},\ {\widehat{z}}_{c} \right) = \arg{\min{\sum_{i = 1}^{N}{\rho\left( f_{i}\left( y_{c},\ z_{c} \right) \right)}}}\
  \label{eq:08}
\end{equation}


Where:


\begin{equation}
  \rho(f) = \gamma^{2}\left( \sqrt{{1 + \left( \frac{f}{\gamma} \right)}^{2}} - 1 \right)\
  \label{eq:09}
\end{equation}


And \(\gamma\) is the soft-\(L_{1}\)scale (0.01m). Fixing R and using
the robust loss makes the centre estimate insensitive both to the
missing invert and to non-lining outliers.

\subsubsection{Polar reduction and invert-gap handling}\label{polar-reduction-and-invert-gap-handling}
\addcontentsline{toc}{paragraph}{Polar reduction and invert-gap
handling}

Each ring is expressed in polar form about its centre using \(\phi_{i}\)
and \(\rho_{i}\):

\begin{equation}
  \phi_{i} = {atan}\left( {z'}_{i} - {\widehat{z}}_{c},\ {y'}_{i} - {\widehat{y}}_{c} \right),\ \rho\sqrt{\left( \ {y'}_{i} - {\widehat{y}}_{c} \right)^{2} + \left( \ {z'}_{i} - {\widehat{z}}_{c} \right)^{2}}\
  \label{eq:10}
\end{equation}

The removed invert is assumed to be the largest empty angular gap, thus,
binning \(\phi\) at \(5^{{^\circ}}\) and locating the longest run of
empty bins gives the gap interval with centre (\(\phi_{g}^{c}\)).

\begin{equation}
  \phi_{i} \leftarrow \ \left( \left( \phi_{i} - \phi_{shift} \right) + \pi \right)mod\ 2\pi + \pi\
  \label{eq:11}
\end{equation}

And \(\phi_{shift} = \phi_{g}^{c} + \pi\)

So that the continuous lining arc occupies a contiguous interval
\(\left\lbrack \phi_{\min},\phi_{\max} \right\rbrack\)in the shifted
frame, with \(\phi = 0\) at the crown

\subsubsection{Longitudinal-joint localisation and segment instancing} \label{longitudinal-joint-localisation-and-segment-instancing}
\addcontentsline{toc}{paragraph}{Longitudinal-joint localisation and
segment instancing}

Longitudinal joints separate the segment instances within a ring. A
longitudinal joint is narrow in \(\phi\) but extends across nearly the
full ring width in \(\omega\), whereas circumferential features are the
opposite; this anisotropy is used to localise joints from the
\emph{Joints point set}. The joint points are transformed to the same
UCS and shifted frame, binned in \(\phi\) (bin width
\({1.5}^{{^\circ}}\)), and each bin \(k\) is scored by its count
\(n_{k}\) weighted by its normalised axial extent,


\begin{equation}
  s_{k} = n_{k}.\frac{\omega_{k,Max} - \omega_{k,Min}}{w_{Total}}\
  \label{eq:12}
\end{equation}


with axial-extent floor \(\tau_{\omega} = 0.70\).

joint angles \(\left( \phi_{j} \right)\)are taken as the peaks of
\(S_{k}\)that meet a relative height threshold
(\(0.35\text{}S_{\max}\)), a relative prominence threshold
(\(0.20\text{}S_{\max}\)) and a minimum separation (\(15^{\circ}\)).
The number of retained peaks is capped at \(n_{seg} - 1\), where
\(n_{seg}\) is the known maximum number of segments visible in the arc.

The joint angles partition the arc into segment instances: a segment
point with an angle \(\phi\) is assigned to the instance bounded by
consecutive joint angles. Instances with fewer than \(5(d + 2)\) points
or an angular span below \(15^{\circ}\)are rejected and recorded.

\subsubsection{Multi-zone polynomial reconstruction of the radial profile}\label{multi-zone-polynomial-reconstruction-of-the-radial-profile}
\addcontentsline{toc}{paragraph}{Multi-zone polynomial reconstruction of
the radial profile}

The radial profile \(\rho(\phi)\) of each segment instance is
reconstructed by a global low-order polynomial augmented with local
``sensitive-zone'' corrections at the joint ends.

Over the interval arc, a polynomial of degree \(d\) (optimised as =4),

\begin{equation}
  \widehat{\rho}_{0}(\phi) = \sum_{m = 0}^{d}{a_{m}\phi^{m}}
  \label{eq:13}
\end{equation}

is fitted by minimising a robust loss with \(\mathcal{l}_{2}\)(ridge)
regularisation on the non-constant coefficients (with ridge weight
\(\beta = 10^{- 3}\)):

\begin{equation}
  \widehat{a} = \arg{\min{\sum_{i = 0}^{a}\rho}}\left( \rho_{i} - {\widehat{\rho}}_{0}\left( \phi_{i} \right) \right) + \beta^{2}\sum_{m = 1}^{d}a_{m}^{2}\
  \label{eq:14}
\end{equation}


\textbf{Sensitive-zone terms.} Within a band
\(\Delta\phi_{s} = 3^{\circ}\)of each end that abuts a joint, a local
correction polynomial of degree \(d_{s} = 3\)is added on top of the
global fit. For the low (\(\phi_{\min}\)) and high (\(\phi_{\max}\))
ends,

\begin{equation}
  \delta_{1}(\phi) = \sum_{l = 1}^{d_{s}}{a_{l}^{(2)}\left( \phi - \phi_{\min} \right)^{l}}\
  for
  \phi \leq \phi_{\min} + \mathrm{\Delta}\phi_{s}\
  \label{eq:15}
\end{equation}


\begin{equation}
  \delta_{2}(\phi) = \sum_{l = 1}^{d_{s}}{a_{l}^{(2)}\left( \phi - \phi_{\max} \right)^{l}}\ for
  \phi \geq \phi_{\max} - \mathrm{\Delta}\phi_{s}\
  \label{eq:16}
\end{equation}



each fitted (same robust + ridge objective) to the residual
\(\rho_{i} - {\widehat{\rho}}_{0}(\phi_{i})\)inside its band.

The per-segment fit quality is summarised by

\begin{equation}
  RMSE = \sqrt{\frac{1}{N_{s}}\sum_{i = 0}^{N_{s}}\left( \rho_{i} - \widehat{\rho}\left( \phi_{i} \right) \right)^{2}}\
  \label{eq:17}
\end{equation}


\subsubsection{Deformation indicators}\label{deformation-indicators}
\addcontentsline{toc}{paragraph}{Deformation indicators}

\emph{Radial deviation:} The signed radial deviation of the lining from
the design circle at angle \(\phi\)is:

\begin{equation}
  \mathrm{\Delta}r(\phi) = \widehat{\rho}(\phi) - R\
  \label{eq:18}
\end{equation}

\emph{Joint dislocation (radial faulting):} At an internal joint between
adjacent instances \(k\) and \(k + 1\), located at
\(\phi_{j} = \frac{1}{2}(\phi_{\max}^{} + \phi_{\min}^{})\), The radial
step (faulting) is the difference of the two reconstructed profiles
evaluated at the joint:

\begin{equation}
  \mathrm{\Delta}_{disl} = {\widehat{\rho}}_{k}\left( \phi_{j} \right) - {\widehat{\rho}}_{k + 1}\left( \phi_{j} \right)\
  \label{eq:19}
\end{equation}

\emph{Joint rotation:} The relative rotation of the two segment faces
meeting at the joint is the change in the angle between the radial
direction and the local surface tangent. With the one-sided slopes
\(\rho_{k}^{'},\rho_{k + 1}^{'} = d\widehat{\rho}/d\phi\)evaluated at
\(\phi_{j}\),

\begin{equation}
  \psi_{k} = {atan}{2\ {\widehat{\rho}}_{k}\left( \phi_{j} \right),\ {\rho'}_{k}}\
  \label{eq:20}
\end{equation}

\emph{Ovalization:} To characterise the cross-sectional
out-of-roundness, an ellipse is fitted to the reconstructed profile
sampled densely as Cartesian points
\(\left( u,v) = (\widehat{\rho}\cos\phi,\widehat{\rho}\sin\phi \right)\).
The general conic

\begin{equation}
  Au^{2} + Buv + Cv^{2} + Du + Ev + F = 0,\ \ \ \ \ \ \ \ \ \ F = - 1,\
  \label{eq:21}
\end{equation}

is fitted in the least-squares sense; the solution is an ellipse when
the discriminant \(\Delta = B^{2} - 4AC < 0\). The semi-axes follow from
the conic coefficients,

\begin{equation}
  a_{e},b_{e} = \sqrt{\frac{- 2\left( {AE}^{2} + {CD}^{2} - BDE + \mathrm{\Delta}F \right)}{\mathrm{\Delta}\left\lbrack (A + C) \mp \sqrt{(A - C)^{2} + B^{2}} \right\rbrack}},\ a_{e} > b_{e}\
  \label{eq:22}
\end{equation}

with major-axis orientation
\(\varphi_{e} = \frac{1}{2}atan2(B,\text{}A - C)\)(resolved to the
major axis). The out-of-roundness (ovality) is

\begin{equation}
  O = a_{e} - b_{e},\ \ \ ellipticity\ \varepsilon = \frac{a_{e} - b_{e}}{R} \times 10^{3}\
  \label{eq:23}
\end{equation}

Serviceability convergence index: The ellipse parameters yield three
distinct, all-valid scalar descriptors of cross-sectional distortion,
and it is essential to distinguish them because they differ by an order
of magnitude and the maintenance standard grades only one of them.

\begin{itemize}
\item
  Out-of-roundness / ellipticity \(\varepsilon = (a_{e} - b_{e})/R\)---
  the total distortion of the section, equal to
  \((D_{\max} - D_{\min})/D_{0}\)with \(D_{\max} = 2a_{e}\),
  \(D_{\min} = 2b_{e}\).
\item
  Major-axis diametric change
  \(\text{\,}C_{maj} = 2(a_{e} - R) = D_{\max} - D_{0}\)--- the outward
  expansion of the longest diameter.
\item
  Clearance convergence \(C_{clr}\)- the reduction of the smallest
  diameter relative to design, i.e. the loss of horizontal clearance.
\end{itemize}

Some Standards ~\citep{MOHURD2018} grades segment convergence as the
\emph{clearance reduction}, not the total out-of-roundness. The
serviceability convergence index used for grading is therefore defined
from the minor diameter,

\[D_{\min} = 2b_{e},\ \ C_{clearance} = \max{\left( \frac{D_{0} - D_{\min}}{D_{0}},0 \right) \times 10^{3}} = \max{\left( \frac{R - b_{e}}{R},0 \right) \times 10^{3}}\]

clamped at zero so that an over-sized (net-outward) ring registers no
clearance loss rather than a spurious negative value. When the
ovalization is not aligned with the principal cross-sectional
directions, the direction-resolved clearance is obtained from the
ellipse radius function

\begin{equation}
  C_{h} = \frac{D_{0} - 2r\left( \theta_{h} \right)}{D_{0}} \times 10^{3}\
  \label{eq:24}
\end{equation}

\begin{equation}
  r(\theta) = \frac{a_{e} - b_{e}}{\sqrt{b_{e}^{2}\cos^{2}\left( \theta - \varphi_{e} \right) + b_{e}^{2}\sin^{2}\left( \theta - \varphi_{e} \right)}}\
  \label{eq:25}
\end{equation}

evaluated at the horizontal direction \(\theta_{h}\); for sections whose
minor axis is near-horizontal, \(C_{h} \approx C_{clr}\), which is used
as the consistency check.

\subsection{Ontology-Based Decision and Information
Layer}\label{ontology-based-decision-and-information-layer}

\subsubsection{The FMEA Reasoning Chain}\label{the-fmea-reasoning-chain}
\addcontentsline{toc}{paragraph}{The FMEA Reasoning Chain}

The reasoning stage of the decision layer includes a four-level chain to
present a degradation mechanism using a measurable indication with a
numeric threshold to come up with a reasonable maintenance action:

\begin{equation}
  FMEA:m_{mechanism} \rightarrow i_{indicator} \rightarrow t_{threshold} \rightarrow \nu_{intervention}\
  \label{eq:26}
\end{equation}

with the four levels indexed by \(L = \left\{ 1,\ \ldots,4 \right\}\).
The mechanism level (L1) names the degradation process; the indicator
level (L2) is the measurable quantity that evidences it; the threshold
level (L3) is the graded band set that maps the indicator to a
structural-health degree; and the intervention level (L4) is the
prescribed maintenance action for that degree.

For a segmental shield metro tunnel, the mechanisms graded in this study
are clearance convergence and joint dislocation (Section \ref{multi-zone-polynomial}) and water
leakage (Section \ref{sonata}). The threshold level is a set of governing
standard tables, causing to produce a health-degree band. Each band is
tagged with its health degree, its intervention tier, and its
originating clause. The authoritative grading scale throughout is the
five-level structural-health degree of the Chinese urban-rail
maintenance standard CJJ/T 289-2018 (Table \ref{tab:02}) ~\citep{MOHURD2018}.

\begin{longtable}[]{@{}
  >{\raggedright\arraybackslash}p{(\columnwidth - 4\tabcolsep) * \real{0.0896}}
  >{\raggedright\arraybackslash}p{(\columnwidth - 4\tabcolsep) * \real{0.7082}}
  >{\raggedright\arraybackslash}p{(\columnwidth - 4\tabcolsep) * \real{0.2123}}@{}}
\caption{Structural-health degree and the corresponding maintenance measure ~\citep{MOHURD2018}}
\label{tab:02}\\ 
\toprule\noalign{}
\begin{minipage}[b]{\linewidth}\raggedright
\textbf{Health degree}
\end{minipage} & \begin{minipage}[b]{\linewidth}\raggedright
\textbf{Maintenance measure (standard)}
\end{minipage} & \begin{minipage}[b]{\linewidth}\raggedright
\textbf{Intervention tier (this study)}
\end{minipage} \\
\midrule\noalign{}
\endhead
\bottomrule\noalign{}
\endlastfoot
Level 1 & Normal upkeep & Acceptable \\
Level 2 & Normal upkeep; attended at the next inspection & Monitor \\
Level 3 & Special monitoring as needed; repair decided from monitoring &
Investigation \\
Level 4 & Restrict use as needed; repair as soon as possible; special
monitoring & Action required \\
Level 5 & Immediately restrict use; repair or replace; special
monitoring & Emergency \\
\end{longtable}

\subsubsection{Source Standards}\label{source-standards}
\addcontentsline{toc}{paragraph}{Source Standards}

The rule base is populated from the prescriptive maintenance and
acceptance standards governing Chinese urban-rail shield tunnels. CJJ/T
289-2018 supplies the structural-health-degree scale (Table \ref{tab:01}1), the
segment-convergence grading bands (Table \ref{tab:03}), and the joint-faulting
grading bands (Table \ref{tab:04}). GB 50208 supplies the waterproofing-grade
criteria used for the leakage assessment (Table \ref{tab:05}).

\begin{longtable}[]{@{}
  >{\raggedright\arraybackslash}p{(\columnwidth - 10\tabcolsep) * \real{0.3197}}
  >{\raggedright\arraybackslash}p{(\columnwidth - 10\tabcolsep) * \real{0.1360}}
  >{\raggedright\arraybackslash}p{(\columnwidth - 10\tabcolsep) * \real{0.1360}}
  >{\raggedright\arraybackslash}p{(\columnwidth - 10\tabcolsep) * \real{0.1360}}
  >{\raggedright\arraybackslash}p{(\columnwidth - 10\tabcolsep) * \real{0.1360}}
  >{\raggedright\arraybackslash}p{(\columnwidth - 10\tabcolsep) * \real{0.1361}}@{}}
\caption{Segment-convergence grading bands ~\citep{MOHURD2018}}
\label{tab:03}\\
\toprule\noalign{}
\begin{minipage}[b]{\linewidth}\raggedright
\textbf{Indicator}
\end{minipage} & \begin{minipage}[b]{\linewidth}\raggedright
\textbf{Level 1}
\end{minipage} & \begin{minipage}[b]{\linewidth}\raggedright
\textbf{Level 2}
\end{minipage} & \begin{minipage}[b]{\linewidth}\raggedright
\textbf{Level 3}
\end{minipage} & \begin{minipage}[b]{\linewidth}\raggedright
\textbf{Level 4}
\end{minipage} & \begin{minipage}[b]{\linewidth}\raggedright
\textbf{Level 5}
\end{minipage} \\
\midrule\noalign{}
\endhead
\bottomrule\noalign{}
\endlastfoot
Straight-joint convergence & $0 \le c \le 5$ & $5 < c \le 8$ & $8 < c \le 12$ & $12 < c \le 16$ & $c > 16$ \\
Staggered-joint convergence & $0 \le c \le 4$ & $4 < c \le 6$ & $6 < c \le 9$ & $9 < c \le 12$ & $c > 12$ \\
\end{longtable}


\begin{longtable}[]{@{}
  >{\raggedright\arraybackslash}p{(\columnwidth - 10\tabcolsep) * \real{0.3199}}
  >{\raggedright\arraybackslash}p{(\columnwidth - 10\tabcolsep) * \real{0.1360}}
  >{\raggedright\arraybackslash}p{(\columnwidth - 10\tabcolsep) * \real{0.1360}}
  >{\raggedright\arraybackslash}p{(\columnwidth - 10\tabcolsep) * \real{0.1360}}
  >{\raggedright\arraybackslash}p{(\columnwidth - 10\tabcolsep) * \real{0.1370}}
  >{\raggedright\arraybackslash}p{(\columnwidth - 10\tabcolsep) * \real{0.1361}}@{}}
\caption{Joint-faulting grading bands ~\citep{MOHURD2018}}
\label{tab:04}\\
\toprule\noalign{}
\begin{minipage}[b]{\linewidth}\raggedright
\textbf{Indicator}
\end{minipage} & \begin{minipage}[b]{\linewidth}\raggedright
\textbf{Level 1}
\end{minipage} & \begin{minipage}[b]{\linewidth}\raggedright
\textbf{Level 2}
\end{minipage} & \begin{minipage}[b]{\linewidth}\raggedright
\textbf{Level 3}
\end{minipage} & \begin{minipage}[b]{\linewidth}\raggedright
\textbf{Level 4}
\end{minipage} & \begin{minipage}[b]{\linewidth}\raggedright
\textbf{Level 5}
\end{minipage} \\
\midrule\noalign{}
\endhead
\bottomrule\noalign{}
\endlastfoot
Joint faulting $f$ (mm) & $0 \le f < 4$ & $4 \le f < 8$ & $8 \le f < 10$ & $10 \le f < 12$ & $f \ge 12$ \\
\end{longtable}

\begin{longtable}[]{@{}
  >{\raggedright\arraybackslash}p{(\columnwidth - 2\tabcolsep) * \real{0.2008}}
  >{\raggedright\arraybackslash}p{(\columnwidth - 2\tabcolsep) * \real{0.8002}}@{}}
\caption{Waterproofing-grade criteria used for the leakage
assessment ~\citep{AQSIQ2011}}
\label{tab:05}\\
\toprule\noalign{}
\begin{minipage}[b]{\linewidth}\raggedright
\textbf{Waterproof grade}
\end{minipage} & \begin{minipage}[b]{\linewidth}\raggedright
\textbf{Leakage Descriptor}
\end{minipage} \\
\midrule\noalign{}
\endhead
\bottomrule\noalign{}
\endlastfoot
Grade 1 & No seepage; no wet stains on the structure surface. \\
Grade 2 & No leakage; minor wet stains permitted -- wet-stain area $\le$
2/1000 of the waterproof area, $\le$ 3 stains per 100 m\textsuperscript{2}, single stain $\le$ 0.2
m\textsuperscript{2}; mean seepage $\le$ 0.05 L/(m\textsuperscript{2}$\cdot$d). \\
Grade 3 & A few leakage points; no linear flow or mud/sand leakage -- $\le$ 7
leakage or wet-stain points per 100 m\textsuperscript{2}, single point $\le$ 2.5 L/d, single
wet stain $\le$ 0.3 m\textsuperscript{2}. \\
Grade 4 & Leakage points present; no linear flow or mud/sand leakage --
mean leakage $\le$ 2 L/(m\textsuperscript{2}$\cdot$d), $\le$ 4 L/(m\textsuperscript{2}$\cdot$d) per 100 m\textsuperscript{2}. \\
\end{longtable}

\subsubsection{LLM-assisted, schema-guided extraction}\label{llm-assisted-schema-guided-extraction}
\addcontentsline{toc}{paragraph}{LLM-assisted, schema-guided extraction}

The rule base is constructed by schema-guided extraction. The large
language model used for this purpose is constrained to populate the
predefined four-field structure of Equation \ref{eq:26} rather than to generate
free-form text. Schema-guided prompting enforces terminological
consistency across standards and yields output that is directly
compatible with rule-based population. Each standard section is
processed in a single long-context pass, so that conditional logic and
cross-references spanning a table and its surrounding clauses are
preserved.

A frontier long-context model was chosen as the primary extraction
engine due to its reasoning ability over dense technical prose and its
robustness to schema-constrained output generation. Decoding is
deterministic (temperature = 0). Categorical-field controlled
vocabularies fix the mechanism (L1) labels so that chains extracted from
different standards are mutually comparable and are refined iteratively
as new terms are encountered. The field definitions and controlled
vocabularies adopted for the metro-shield instantiation are summarised
in Table \ref{tab:06}.

\begin{longtable}[]{@{}
  >{\raggedright\arraybackslash}p{(\columnwidth - 6\tabcolsep) * \real{0.0897}}
  >{\raggedright\arraybackslash}p{(\columnwidth - 6\tabcolsep) * \real{0.1323}}
  >{\raggedright\arraybackslash}p{(\columnwidth - 6\tabcolsep) * \real{0.4960}}
  >{\raggedright\arraybackslash}p{(\columnwidth - 6\tabcolsep) * \real{0.2820}}@{}}
\caption{Four-level chain: field definitions and controlled
vocabularies for the metro-shield instantiation.}
\label{tab:06}\\ 
\toprule\noalign{}
\begin{minipage}[b]{\linewidth}\raggedright
\textbf{FMEA Level}
\end{minipage} & \begin{minipage}[b]{\linewidth}\raggedright
\textbf{Field}
\end{minipage} & \begin{minipage}[b]{\linewidth}\raggedright
\textbf{Controlled vocabulary/ form}
\end{minipage} & \begin{minipage}[b]{\linewidth}\raggedright
\textbf{Example (metro shield)}
\end{minipage} \\
\midrule\noalign{}
\endhead
\bottomrule\noalign{}
\endlastfoot
L1 & Mechanism & Convergence deformation, Joint dislocation, Water
leakage, Joint opening, Cracking, Spalling & Clearance convergence of
the lining ring \\
L2 & Indicator & Convergence, dislocation/faulting (mm), wet area (m²)
and leakage class, opening (mm) & Horizontal clearance convergence, of
diameter \\
L3 & Threshold & Graded band set on the five-level CJJ/T 289 health
degree, with source clause & 0-4 / 4-6 / 6-9 / 9-12 / \textgreater12 \\
L4 & Intervention & Routine upkeep, Monitoring, Investigation, Repair,
Reinforcement, Replacement & Reinforcement / clearance-restoration
grouting \\
\end{longtable}

\subsubsection{Population of the present tunnel and modality
mapping}\label{population-of-the-present-tunnel-and-modality-mapping}
\addcontentsline{toc}{paragraph}{Population of the present tunnel and
modality mapping}

For the present tunnel, the indicator level (L2) of each chain is
generated with instance-level condition data from two evidence channels.
The geometric channel supplies the structural indicators from the
deformation pipeline: the horizontal clearance convergence and the
per-ring maximum joint dislocation enter the convergence and dislocation
chains, respectively. The image channel supplies the leakage indicator:
a vision--language model reads the segmented leakage map and the
radial-deviation heat-map of each axial window, together with the
corresponding per-ring numeric table, and returns a structured record
giving the leakage class (none, wet-stain, seepage, dripping,
linear-leakage, gushing), the total wet area, and the located leakage
patches.

\subsubsection{Rule-driven grading and
intervention}\label{rule-driven-grading-and-intervention}
\addcontentsline{toc}{paragraph}{Rule-driven grading and intervention}

The decision layer uses the extracted rule base as the grader and
iterates the four-level chain for each ring. For a measured mechanism
(L1), the governing chain is chosen by the mechanism label, the
indicator unit, and the joint type of the ring. The measured indicator
(L2) is compared with the graded threshold bands (L3) of that chain, and
the health degree is the band in which the indicator lies; an indicator
value outside the outermost band is assigned to the nearest extreme
band. The structural condition of the ring is the most severe category
within the graded mechanisms. The leakage condition is derived by
comparing the vision-derived leakage class and, where applicable, the
wet area against the waterproofing-grade criteria of Table 4. The
overall condition of a tunnel section is the worst of its worst
ring-level structural degree, and its leakage degree, and its prescribed
maintenance action (L4) is the intervention of the chain that determines
the governing degree, reported with its tier (Table \ref{Tab:01}) and source
clause.

\subsubsection{Scope of the decision
layer}\label{scope-of-the-decision-layer}
\addcontentsline{toc}{paragraph}{Scope of the decision layer}

The decision layer scores for the four levels of Eq. (\ref{eq:26}). The structural
grade (convergence and dislocation) is obtained from the numeric
standard bands of Tables 2 and 3. It is offset-immune due to the
clearance-convergence formulation. Leakage is classified qualitatively
by visual inspection and related to the descriptor criteria in Table \ref{tab:04}.
The area criteria are calibrated from the acceptance standard and
engineering practice, so leakage-governed results are reported as
subject to field verification.

\section{Results and discussion}\label{results-and-discussion}

\subsection{Leakage Prediction}\label{leakage-prediction-1}

The leakage-segmentation experiments were designed to follow a single
question across a grid of conditions: how much of the model's accuracy
survives while reducing the amount of labelled data for the training
section. Three labelled-data budgets (25\%, 50\%, and 100\% of the
training files) were crossed with two transfer regimes, full fine-tuning
of the encoder and linear probing of a frozen encoder, and the whole
grid was run twice: once on the S3DIS leakage data alone and once with
the Australian Reference Tunnel (ART) added to the training pool.

The full-fine-tuning results are presented inTable \ref{tab:07}. It can be
seen that the model trained on S3DIS alone rises from an accuracy of
0.93 and a macro-F1 of 0.93 at the quarter-data budget to 0.97 on both
measures at full supervision, with Cohen's $\kappa$ and MCC moving from
0.89--0.90 to 0.95. The linear-probing columns of the same dataset begin
much lower and climb more steeply, from an accuracy of 0.69 and a $\kappa$ of
0.58 at 25\% to 0.93 and 0.89 at full data, so that the separation
between the two regimes is widest at the smallest budget and narrows as
labels are added. Adding ART changes the picture chiefly at the weak end
of the grid: the linear-probe configuration at 25\% rises from an
accuracy of 0.69 to 0.83 and its $\kappa$ from 0.58 to 0.75, whereas the
full-fine-tuning configurations, already high on S3DIS alone, shift only
within a few hundredths and at full data sit marginally below their
S3DIS-only counterparts at 0.96 against 0.97 accuracy.


\begin{longtable}[]{@{}
  >{\raggedright\arraybackslash}p{(\columnwidth - 12\tabcolsep) * \real{0.2431}}
  >{\raggedright\arraybackslash}p{(\columnwidth - 12\tabcolsep) * \real{0.1261}}
  >{\raggedright\arraybackslash}p{(\columnwidth - 12\tabcolsep) * \real{0.1261}}
  >{\raggedright\arraybackslash}p{(\columnwidth - 12\tabcolsep) * \real{0.1261}}
  >{\raggedright\arraybackslash}p{(\columnwidth - 12\tabcolsep) * \real{0.1261}}
  >{\raggedright\arraybackslash}p{(\columnwidth - 12\tabcolsep) * \real{0.1261}}
  >{\raggedright\arraybackslash}p{(\columnwidth - 12\tabcolsep) * \real{0.1261}}@{}}
\caption{performance evaluation across different configurations.}
\label{tab:07}\\ 
\toprule\noalign{}
\begin{minipage}[b]{\linewidth}\raggedright
\end{minipage} &
\multicolumn{6}{>{\raggedright\arraybackslash}p{(\columnwidth - 12\tabcolsep) * \real{0.7569} + 10\tabcolsep}@{}}{%
\begin{minipage}[b]{\linewidth}\raggedright
\emph{\textbf{S3DIS+ ART}}
\end{minipage}} \\
\midrule\noalign{}
\endhead
\bottomrule\noalign{}
\endlastfoot
&
\multicolumn{3}{>{\raggedright\arraybackslash}p{(\columnwidth - 12\tabcolsep) * \real{0.3784} + 4\tabcolsep}}{%
\textbf{Full}} &
\multicolumn{3}{>{\raggedright\arraybackslash}p{(\columnwidth - 12\tabcolsep) * \real{0.3784} + 4\tabcolsep}@{}}{%
\textbf{Linear Probe}} \\
\textbf{Database size (\%):} & 25 & 50 & 100 & 25 & 50 & 100 \\
\textbf{Accuracy} & 0.95 & 0.95 & 0.96 & 0.83 & 0.90 & 0.90 \\
\textbf{F1 (macro)} & 0.95 & 0.95 & 0.95 & 0.84 & 0.90 & 0.91 \\
\textbf{F1 (Weighted)} & 0.95 & 0.96 & 0.96 & 0.84 & 0.90 & 0.91 \\
\textbf{Cohen's Kappa} & 0.92 & 0.93 & 0.93 & 0.75 & 0.84 & 0.86 \\
\textbf{MCC} & 0.93 & 0.93 & 0.93 & 0.78 & 0.86 & 0.87 \\
&
\multicolumn{6}{>{\raggedright\arraybackslash}p{(\columnwidth - 12\tabcolsep) * \real{0.7569} + 10\tabcolsep}@{}}{%
\emph{\textbf{S3DIS}}} \\
&
\multicolumn{3}{>{\raggedright\arraybackslash}p{(\columnwidth - 12\tabcolsep) * \real{0.3784} + 4\tabcolsep}}{%
\textbf{Full}} &
\multicolumn{3}{>{\raggedright\arraybackslash}p{(\columnwidth - 12\tabcolsep) * \real{0.3784} + 4\tabcolsep}@{}}{%
\textbf{Linear Probe}} \\
\textbf{Database size (\%):} & 25 & 50 & 100 & 25 & 50 & 100 \\
\textbf{Accuracy} & 0.93 & 0.96 & 0.97 & 0.69 & 0.82 & 0.93 \\
\textbf{F1 (macro)} & 0.93 & 0.96 & 0.97 & 0.74 & 0.83 & 0.93 \\
\textbf{F1 (Weighted)} & 0.93 & 0.96 & 0.97 & 0.69 & 0.82 & 0.93 \\
\textbf{Cohen's Kappa} & 0.89 & 0.93 & 0.95 & 0.58 & 0.73 & 0.89 \\
\textbf{MCC} & 0.90 & 0.94 & 0.95 & 0.65 & 0.77 & 0.90 \\
\end{longtable}

The predicted segmentations besides the ground-truth labelling for all
twelve configurations are presented in Figure \ref{fig:Fig5}. The qualitative trend
follows the tabulated one: the full-fine-tuning panels reproduce the
position and extent of the blue leakage patches closely at every budget,
and the differences among 25\%, 50\%, and 100\% are confined to the
sharpness of patch boundaries and a few small, isolated regions. The
linear-probe panels show the data-budget effect more plainly, with the
25\% maps undermarking the smaller leakage areas and recovering
progressively more of them at 50\% and 100\%.


\begin{figure}[htbp]
  \centering
  \includegraphics[width=0.92\textwidth]{Fig_PDF/Picture052.pdf}
  \caption{ a comparison between the ground truth and different
configurations}
  \label{fig:Fig5}
\end{figure}
The given configurations are then compared against nine point-cloud
baselines reported by  ~\citep{Jiang2025a,Jiang2025b} on the same leakage
benchmark in Table \ref{tab:08}. The highest leakage IoU in the table (94.34)
belongs to the full-fine-tuning configuration trained on S3DIS at full
supervision; the next values beneath it are 92.37 and 92.17 from the
50\%-S3DIS and full-data S3DIS+ART full-fine-tuning runs, and the
strongest external baseline, the BBE-Net ensemble, records 88.95. On
mean accuracy, the same full-data S3DIS configuration reaches 98.60,
above the 96.01 of BBE-Net, while on mean IoU, the 94.26 of BBE-Net sits
just above the 93.90 of that configuration.

The internal ordering within the present study repeats the regime and
budget pattern of Table \ref{tab:07} and expresses the dataset contrast in IoU
terms. Full fine-tuning exceeds linear probing at every matched setting.
The ART contribution again concentrates at the weak end of the grid: the
25\% linear-probe leakage IoU rises from 45.41 on S3DIS to 70.02 with
ART, whereas at full-data full fine-tuning the S3DIS-only value of 94.34
stands above the 92.17 obtained with ART.

\begin{longtable}[]{@{}
  >{\raggedright\arraybackslash}p{(\columnwidth - 10\tabcolsep) * \real{0.2085}}
  >{\raggedright\arraybackslash}p{(\columnwidth - 10\tabcolsep) * \real{0.2361}}
  >{\raggedright\arraybackslash}p{(\columnwidth - 10\tabcolsep) * \real{0.3057}}
  >{\raggedright\arraybackslash}p{(\columnwidth - 10\tabcolsep) * \real{0.1015}}
  >{\raggedright\arraybackslash}p{(\columnwidth - 10\tabcolsep) * \real{0.0741}}
  >{\raggedright\arraybackslash}p{(\columnwidth - 10\tabcolsep) * \real{0.0741}}@{}}
\caption{Results comparison with the previous studies.}
\label{tab:08}\\ 
\toprule\noalign{}
\begin{minipage}[b]{\linewidth}\raggedright
\textbf{Source}
\end{minipage} & \begin{minipage}[b]{\linewidth}\raggedright
\textbf{method}
\end{minipage} & \begin{minipage}[b]{\linewidth}\raggedright
\textbf{Backbone family}
\end{minipage} & \begin{minipage}[b]{\linewidth}\raggedright
\textbf{L-IoU}
\end{minipage} & \begin{minipage}[b]{\linewidth}\raggedright
\textbf{mIoU}
\end{minipage} & \begin{minipage}[b]{\linewidth}\raggedright
\textbf{mAcc}
\end{minipage} \\
\midrule\noalign{}
\endhead
\bottomrule\noalign{}
\endlastfoot
\multirow{12}{=}{This study} & \multirow{12}{=}{PTv3 backbone} & Sonata
(Full, S3DIS+ART, 25\%) & 88.52 & 90.70 & 97.80 \\
& & Sonata (Full, S3DIS+ART, 50\%) & 91.31 & 88.10 & 93.50 \\
& & Sonata (Full, S3DIS+ART, 100\%) & 92.17 & 91.50 & 98.00 \\
& & Sonata (Lin., S3DIS+ART, 25\%) & 70.02 & 74.20 & 92.10 \\
& & Sonata (Lin., S3DIS+ART, 50\%) & 81.79 & 82.50 & 95.30 \\
& & Sonata (Lin., S3DIS+ART, 100\%) & 83.19 & 83.80 & 95.70 \\
& & Sonata (Full, S3DIS, 25\%) & 87.59 & 87.40 & 96.90 \\
& & Sonata (Full, S3DIS, 50\%) & 92.37 & 91.70 & 98.10 \\
& & Sonata (Full, S3DIS, 100\%) & 94.34 & 93.90 & 98.60 \\
& & Sonata (Lin., S3DIS, 25\%) & 45.41 & 60.90 & 85.60 \\
& & Sonata (Lin., S3DIS, 50\%) & 67.65 & 73.10 & 91.40 \\
& & Sonata (Lin., S3DIS, 100\%) & 88.01 & 87.30 & 96.80 \\
\multirow{2}{=}{(Jiang et al., 2025a)} & GEPT-Net & Point Transformer v3
+ FPFH & 85.31 & 92.38 & 94.50 \\
& PTv2 & Point Transformer v2 & 75.17 & 87.17 & 90.24 \\
\multirow{7}{=}{(Jiang et al., 2025b)} & BBE-Net & Ensemble (SpU + ST +
PT) & 88.95 & 94.26 & 96.01 \\
& RandLA-Net & RandLA & 74.03 & 86.74 & 92.38 \\
& Minkowski U-Net & Sparse conv & 71.57 & 85.35 & 91.63 \\
& Sparse U-Net & Sparse conv & 71.15 & 85.21 & 92.74 \\
& Stratified Transformer & Transformer & 67.73 & 83.44 & 88.07 \\
& Point Transformer v1 & Point Transformer & 66.61 & 82.74 & 88.47 \\
& SPVCNN & Sparse point-voxel & 55.54 & 76.89 & 82.64 \\
\end{longtable}

Figure \ref{fig:Fig6} shows the confusion matrix of the strongest configuration, full
fine-tuning on S3DIS at full supervision, in both count and
row-normalised form. The Joints, Segments, and Pockets classes are
recovered without error at the normalised level, each reading 1.00 on
the diagonal. The leakage row reads 0.94 on the diagonal, with the
remaining true-leakage points distributed as 0.05 into Joints and 0.01
into Pockets and a negligible fraction into Segments; in counts,
2,225,756 leakage points are labelled correctly while 117,056 are
assigned to Joints and 16,307 to Pockets.

Moreover, Figure \ref{fig:Fig7} resolves the same configuration into per-class
precision, recall, F1, and IoU. The leakage class records a precision of
1.000 with a recall of 0.943, so that no point of another class was
assigned to leakage while 5.7\% of true leakage points received other
labels; the Joints class shows the complementary pattern, a recall of
1.000 against a precision of 0.840, in line with the leakage points that
the confusion matrix placed in Joints. The Segments class reads 1.000 on
all four measures, and the Pockets class records a precision of 0.974, a
recall of 1.000, an F1 of 0.987, and an IoU of 0.974; expressed as IoU,
the four classes stand at 0.945 for leakage, 0.840 for joints, 1.000 for
segments, and 0.974 for pockets.

\begin{figure}[H]
  \centering
  \includegraphics[width=0.92\textwidth]{Fig_PDF/Picture06.pdf}
  \caption{ Confusion matrix for \emph{Sonata (Full, S3DIS, 100\%)}}
  \label{fig:Fig6}
\end{figure}

\begin{figure}[H]
  \centering
  \includegraphics[width=0.92\textwidth]{Fig_PDF/Picture07.pdf}
  \caption{ Per-class metrics for \emph{Sonata (Full, S3DIS, 100\%)}}
  \label{fig:Fig7}
\end{figure}


\subsection{Deformation}\label{deformation}
Once the lining had been separated into its structural elements, the
geometry of each ring was reconstructed and reduced to a small set of
deformation indicators, so that the qualitative segmentation of the
previous stage could be read as measured shape. The reconstruction was
applied ring by ring to the two held-out tunnels, the thirteen rings of
Tunnel 8 and the twelve rings of Tunnel 9, and every ring was
characterised by its clearance convergence, its ovalization and
ellipticity, the largest joint dislocation and rotation found along it,
and the residual of the ellipse fit. Table \ref{tab:09} records these quantities
for the 25 rings.

Across both tunnels the clearance convergence remains within a narrow
band, running from 1.59 to 4.22 ‰D in Tunnel 8 and from 2.71 to 4.93 ‰D
in Tunnel 9, while the ovalization spans roughly 40 to 52 mm and the
ellipticity 14.6 to 18.7 ‰. The maximum joint dislocation per ring
reaches 5.4 mm in Tunnel 8 and 5.2 mm in Tunnel 9, and the maximum joint
rotation is more variable, mostly between 2° and 9° but rising to 12.49°
at ring 12 of Tunnel 8 and 12.07° at ring 8 of Tunnel 9. The ellipse-fit
RMSE stays between 4.48 and 7.04 mm across all rings, with the larger
residuals coinciding with the rings that carry the highest ovalization.

\begin{longtable}[]{@{}
  >{\raggedright\arraybackslash}p{(\columnwidth - 14\tabcolsep) * \real{0.0919}}
  >{\raggedright\arraybackslash}p{(\columnwidth - 14\tabcolsep) * \real{0.0654}}
  >{\raggedright\arraybackslash}p{(\columnwidth - 14\tabcolsep) * \real{0.1404}}
  >{\raggedright\arraybackslash}p{(\columnwidth - 14\tabcolsep) * \real{0.1404}}
  >{\raggedright\arraybackslash}p{(\columnwidth - 14\tabcolsep) * \real{0.1404}}
  >{\raggedright\arraybackslash}p{(\columnwidth - 14\tabcolsep) * \real{0.1404}}
  >{\raggedright\arraybackslash}p{(\columnwidth - 14\tabcolsep) * \real{0.1404}}
  >{\raggedright\arraybackslash}p{(\columnwidth - 14\tabcolsep) * \real{0.1404}}@{}}
\caption{Per-ring deformation indicators for the two held-out test
tunnels.}
\label{tab:09}\\
\toprule\noalign{}
\begin{minipage}[b]{\linewidth}\raggedright
\textbf{Tunnel}
\end{minipage} & \begin{minipage}[b]{\linewidth}\raggedright
\textbf{Ring}
\end{minipage} & \begin{minipage}[b]{\linewidth}\raggedright
\textbf{Clearance conv. (\textperthousand D)}
\end{minipage} & \begin{minipage}[b]{\linewidth}\raggedright
\textbf{Ovalization (mm)}
\end{minipage} & \begin{minipage}[b]{\linewidth}\raggedright
\textbf{Ellipticity (\textperthousand)}
\end{minipage} & \begin{minipage}[b]{\linewidth}\raggedright
\textbf{Max joint dislocation (mm)}
\end{minipage} & \begin{minipage}[b]{\linewidth}\raggedright
\textbf{Max joint rotation ($^{\circ}$)}
\end{minipage} & \begin{minipage}[b]{\linewidth}\raggedright
\textbf{Fit RMSE (mm)}
\end{minipage} \\
\midrule\noalign{}
\endhead
\bottomrule\noalign{}
\endlastfoot
8 & 1 & 2.24 & 41.5 & 15.1 & 3 & 5.06 & 4.48 \\
8 & 2 & 2.4 & 44.3 & 16.1 & 5.4 & 5.45 & 4.54 \\
8 & 3 & 2.17 & 43.5 & 15.8 & 1.8 & 5.31 & 5.01 \\
8 & 4 & 1.59 & 41.7 & 15.2 & 2.2 & 3.48 & 5.23 \\
8 & 5 & 3.38 & 49 & 17.8 & 2.7 & 2.48 & 4.92 \\
8 & 6 & 2.94 & 47 & 17.1 & 1.6 & 8.21 & 4.53 \\
8 & 7 & 2.48 & 44 & 16 & 2.6 & 3.04 & 5.53 \\
8 & 8 & 2.29 & 45 & 16.4 & 4.7 & 5.46 & 5.67 \\
8 & 9 & 4.22 & 50.2 & 18.2 & 2.6 & 2.15 & 6.04 \\
8 & 10 & 2.42 & 43.9 & 15.9 & 3.1 & 3.55 & 5.08 \\
8 & 11 & 3.91 & 48.5 & 17.6 & 3.3 & 1.99 & 5.7 \\
8 & 12 & 3.8 & 49.4 & 18 & 4.3 & 12.49 & 7.04 \\
8 & 13 & 3.36 & 47.1 & 17.1 & 3.9 & 4.17 & 6.54 \\
9 & 1 & 4.6 & 50.7 & 18.5 & 2.5 & 5.59 & 5.37 \\
9 & 2 & 3.73 & 46.6 & 16.9 & 3.9 & 4.17 & 5.8 \\
9 & 3 & 4.93 & 51.5 & 18.7 & 2.5 & 4.56 & 5.25 \\
9 & 4 & 4.07 & 48.8 & 17.7 & 3.4 & 8.34 & 5.72 \\
9 & 5 & 3.57 & 46.2 & 16.8 & 5.2 & 5.88 & 5.96 \\
9 & 6 & 4.5 & 49.9 & 18.1 & 3.9 & 9.57 & 6.71 \\
9 & 7 & 3.88 & 46.1 & 16.8 & 1 & 3.91 & 6.05 \\
9 & 8 & 3.2 & 44.3 & 16.1 & 3.5 & 12.07 & 6.86 \\
9 & 9 & 2.71 & 40.1 & 14.6 & 4 & 2.55 & 5.86 \\
9 & 10 & 3.66 & 43.2 & 15.7 & 2.5 & 4.36 & 5.09 \\
9 & 11 & 3.6 & 42.9 & 15.6 & 4.2 & 7.96 & 5.01 \\
9 & 12 & 4.34 & 45.6 & 16.6 & 3 & 7.67 & 5.55 \\
\end{longtable}

Because the ovalization is the indicator most sensitive to the fitting
procedure, it was computed by two independent routes and compared. Table
\ref{tab:10} lists, for every ring, the ovalization obtained from the polynomial
reconstruction and from a raw conic fit, together with their absolute
difference. The two estimates track each other closely: the mean
absolute difference is 1.15 mm for Tunnel 8 and 1.62 mm for Tunnel 9,
with an overall mean of 1.37 mm, and the largest single discrepancy is
3.6 mm, at ring 9 of Tunnel 8, where the polynomial value of 50.18 mm
meets a raw value of 46.58 mm.

\begin{longtable}[]{@{}
  >{\raggedright\arraybackslash}p{(\columnwidth - 8\tabcolsep) * \real{0.0892}}
  >{\raggedright\arraybackslash}p{(\columnwidth - 8\tabcolsep) * \real{0.0701}}
  >{\raggedright\arraybackslash}p{(\columnwidth - 8\tabcolsep) * \real{0.3147}}
  >{\raggedright\arraybackslash}p{(\columnwidth - 8\tabcolsep) * \real{0.3018}}
  >{\raggedright\arraybackslash}p{(\columnwidth - 8\tabcolsep) * \real{0.2242}}@{}}
\caption{Ring-by-ring comparison of ovalization computed from the
polynomial reconstruction and from the raw conic fit, with the absolute
difference per ring and the per-tunnel and overall mean and maximum
differences.}
\label{tab:10}\\
\toprule\noalign{}
\begin{minipage}[b]{\linewidth}\raggedright
\textbf{Tunnel}
\end{minipage} & \begin{minipage}[b]{\linewidth}\raggedright
\textbf{Ring}
\end{minipage} & \begin{minipage}[b]{\linewidth}\raggedright
\textbf{Ovalization, polynomial (mm)}
\end{minipage} & \begin{minipage}[b]{\linewidth}\raggedright
\textbf{Ovalization, raw conic (mm)}
\end{minipage} & \begin{minipage}[b]{\linewidth}\raggedright
\textbf{Abs. difference (mm)}
\end{minipage} \\
\midrule\noalign{}
\endhead
\bottomrule\noalign{}
\endlastfoot
8 & 1 & 41.48 & 40.93 & 0.55 \\
8 & 2 & 44.26 & 43.89 & 0.36 \\
8 & 3 & 43.45 & 43.97 & 0.52 \\
8 & 4 & 41.71 & 42.04 & 0.33 \\
8 & 5 & 49.03 & 48.78 & 0.26 \\
8 & 6 & 47.01 & 45.81 & 1.19 \\
8 & 7 & 43.99 & 42.17 & 1.82 \\
8 & 8 & 45 & 45.67 & 0.67 \\
8 & 9 & 50.18 & 46.58 & 3.6 \\
8 & 10 & 43.85 & 44.36 & 0.51 \\
8 & 11 & 48.51 & 46.29 & 2.22 \\
8 & 12 & 49.41 & 46.87 & 2.54 \\
8 & 13 & 47.12 & 47.45 & 0.34 \\
\textbf{T8} & \emph{mean} & & & 1.15 \\
\textbf{T8} & \emph{max} & ~ & ~ & 3.6 \\
9 & 1 & 50.75 & 48.37 & 2.38 \\
9 & 2 & 46.6 & 45.99 & 0.61 \\
9 & 3 & 51.48 & 48.81 & 2.66 \\
9 & 4 & 48.78 & 46.09 & 2.7 \\
9 & 5 & 46.18 & 46.41 & 0.23 \\
9 & 6 & 49.86 & 47.38 & 2.48 \\
9 & 7 & 46.11 & 45.88 & 0.23 \\
9 & 8 & 44.29 & 42.62 & 1.67 \\
9 & 9 & 40.07 & 40.89 & 0.82 \\
9 & 10 & 43.17 & 40.49 & 2.68 \\
9 & 11 & 42.85 & 42.6 & 0.25 \\
9 & 12 & 45.65 & 42.95 & 2.7 \\
\textbf{T9} & \emph{mean} & & & 1.62 \\
\textbf{T9} & \emph{max} & & & 2.7 \\
\textbf{All} & \emph{mean} & & & 1.37 \\
\textbf{All} & \emph{max} & ~ & ~ & 3.6 \\
\end{longtable}

Figure \ref{fig:Fig8} plots the same comparison, with the polynomial ovalization on
one axis and the raw-conic ovalization on the other. The points for both
tunnels fall along the diagonal, with the Tunnel 8 cloud (panel a, mean
absolute difference 1.1 mm) sitting marginally tighter than the Tunnel 9
cloud (panel b, 1.6 mm), and the few points that depart from the line
are the high-ovalization rings already identified in Table 10.

\begin{figure}[htbp]
  \centering
  \includegraphics[width=0.92\textwidth]{Fig_PDF/Picture08.pdf}
  \caption{ Agreement between the polynomial (pipeline) and raw-conic
ovalization estimates for (a) Tunnel 8 and (b) Tunnel 9}
  \label{fig:Fig8}
\end{figure}

The joint behaviour of each ring is detailed in Table \ref{tab:11}, which
separates the mean and maximum absolute dislocation from the mean and
maximum absolute rotation. The mean absolute dislocation per ring lies
between 0.38 and 2.86 mm and the per-ring maximum between 0.98 and 5.37
mm. It can be seen that the dislocations are small and uniform along
both tunnels. The rotation values are less uniform: mean absolute
rotation per ring ranges from about 1.1° to 6.3°, and the per-ring
maximum reaches 12.49° at ring 12 of Tunnel 8 and 12.07° at ring 8 of
Tunnel \ref{tab:09}, the same two rings that carry the largest rotations in Table
9.

\begin{longtable}[]{@{}
  >{\raggedright\arraybackslash}p{(\columnwidth - 10\tabcolsep) * \real{0.1097}}
  >{\raggedright\arraybackslash}p{(\columnwidth - 10\tabcolsep) * \real{0.1098}}
  >{\raggedright\arraybackslash}p{(\columnwidth - 10\tabcolsep) * \real{0.1982}}
  >{\raggedright\arraybackslash}p{(\columnwidth - 10\tabcolsep) * \real{0.1982}}
  >{\raggedright\arraybackslash}p{(\columnwidth - 10\tabcolsep) * \real{0.1968}}
  >{\raggedright\arraybackslash}p{(\columnwidth - 10\tabcolsep) * \real{0.1873}}@{}}
\caption{Per-ring joint behavior for the two test tunnels: mean
and maximum absolute joint dislocation (mm) and mean and maximum
absolute joint rotation (°).}
\label{tab:11}\\
\toprule\noalign{}
\begin{minipage}[b]{\linewidth}\raggedright
\textbf{Tunnel}
\end{minipage} & \begin{minipage}[b]{\linewidth}\raggedright
\textbf{Ring}
\end{minipage} & \begin{minipage}[b]{\linewidth}\raggedright
\textbf{Mean abs. dislocation (mm)}
\end{minipage} & \begin{minipage}[b]{\linewidth}\raggedright
\textbf{Max abs. dislocation (mm)}
\end{minipage} & \begin{minipage}[b]{\linewidth}\raggedright
\textbf{Mean abs. rotation (°)}
\end{minipage} & \begin{minipage}[b]{\linewidth}\raggedright
\textbf{Max abs. rotation (°)}
\end{minipage} \\
\midrule\noalign{}
\endhead
\bottomrule\noalign{}
\endlastfoot
8 & 1 & 1.25 & 3.01 & 3.26 & 5.06 \\
8 & 2 & 2.49 & 5.37 & 2.81 & 5.45 \\
8 & 3 & 1.27 & 1.81 & 3.31 & 5.31 \\
8 & 4 & 1.2 & 2.22 & 1.43 & 3.48 \\
8 & 5 & 2.12 & 2.68 & 1.21 & 2.48 \\
8 & 6 & 0.67 & 1.59 & 4.14 & 8.21 \\
8 & 7 & 1.43 & 2.64 & 2.31 & 3.04 \\
8 & 8 & 2.36 & 4.7 & 2.18 & 5.46 \\
8 & 9 & 1.7 & 2.62 & 1.43 & 2.15 \\
8 & 10 & 1.92 & 3.14 & 2.61 & 3.55 \\
8 & 11 & 1.53 & 3.26 & 1.08 & 1.99 \\
8 & 12 & 2.68 & 4.28 & 6.27 & 12.49 \\
8 & 13 & 2.63 & 3.85 & 1.99 & 4.17 \\
9 & 1 & 1.73 & 2.47 & 4.24 & 5.59 \\
9 & 2 & 1.92 & 3.94 & 2.39 & 4.17 \\
9 & 3 & 1.42 & 2.47 & 3.09 & 4.56 \\
9 & 4 & 2.18 & 3.44 & 4.17 & 8.34 \\
9 & 5 & 2.41 & 5.15 & 3.35 & 5.88 \\
9 & 6 & 2.86 & 3.89 & 4.42 & 9.57 \\
9 & 7 & 0.38 & 0.98 & 3.16 & 3.91 \\
9 & 8 & 1.97 & 3.46 & 6 & 12.07 \\
9 & 9 & 2.58 & 4.04 & 1.37 & 2.55 \\
9 & 10 & 1.91 & 2.54 & 3.03 & 4.36 \\
9 & 11 & 1.75 & 4.15 & 3.21 & 7.96 \\
9 & 12 & 2.33 & 2.98 & 5.12 & 7.67 \\
\end{longtable}

Figure \ref{fig:Fig9} renders these quantities on the reconstructed three-dimensional
lining. Panel (a) shows the radial-deviation heatmap, in which the
deviation from the nominal design circle varies continuously around and
along the tunnel, the crown region appearing in the negative end of the
scale and the haunch and invert regions in the positive end. Panels (b)
and (c) map the joint dislocation and joint rotation onto the segment
patches, the dislocation field in panel (b) showing pale, evenly
distributed patches and the rotation field in panel (c) showing a
sparser set of darker patches concentrated at particular ring positions,
including the rings that registered the highest rotation magnitudes in
Tables 9 and 11. Also, cross-sections of each ring along the radial
deformation are presented in Figure \ref{fig:FigA1} and Figure \ref{fig:FigA2} in Appendix
\ref{appendix-a--per-ring-deformation-cross-section}.

The per-ring quantities are then serialised for the stage that follows.
For each ring, the clearance convergence and the maximum joint
dislocation are written together with the ovalization, ellipticity,
rotation, and fit RMSE into a per-ring record keyed by tunnel and ring
number; the per-joint dislocation and rotation values of Table 11 and
the ovalization-agreement statistics of Table 10 are stored alongside
them. These records, exported as the JSON and CSV files read by the
ontology-based decision stage of Section 4.3, hold the per-ring
clearance convergence and joint dislocation that form the numeric
channel of that stage.

\begin{figure}[htbp]
  \centering
  \begin{subfigure}{\textwidth}
    \centering
    \includegraphics[width=0.92\textwidth]{Fig_PDF/Picture09a.pdf}
    \caption{Radial-deviation heatmap}
    \label{fig:Fig9a}
  \end{subfigure}
\end{figure}

\begin{figure}[htbp]\ContinuedFloat
  \centering
  \begin{subfigure}{\textwidth}
    \centering
    \includegraphics[width=0.92\textwidth]{Fig_PDF/Picture09b.pdf}
    \caption{Joint dislocation}
    \label{fig:Fig9b}
  \end{subfigure}
\end{figure}

\begin{figure}[htbp]\ContinuedFloat
  \centering
  \begin{subfigure}{\textwidth}
    \centering
    \includegraphics[width=0.92\textwidth]{Fig_PDF/Picture09c.pdf}
    \caption{Joint rotation}
    \label{fig:Fig9c}
  \end{subfigure}
  \caption{Three-dimensional rendering of the reconstructed test-tunnel lining: (a) radial-deviation heatmap, (b) joint dislocation, and (c) joint rotation.}
  \label{fig:Fig9}
\end{figure}



\subsection{Ontology-Based Decision
Support}\label{ontology-based-decision-support}

The schema-guided extraction resulted in 51 FMEA reasoning chains,
derived from the seven sections of the Chinese urban-rail maintenance
codes governing, and Table \ref{tab:12} organizes these chains based on the
degradation mechanism. Eleven classes of mechanisms are shown. The most
populated is water leakage, with ten chains, then convergence
deformation (eight), longitudinal deformation (six). Joint dislocation
and joint opening each account for five chains. The chains are
classified into ten component categories, where segment lining is the
most dominant category with twenty-four chains, followed by structural
connection with eleven chains, while the remaining categories cover
cut-and-cover structures, track bed, waterproofing system, bolts and
ancillary openings.

Each chain has its indicators, the indicator units, and the source
standard, so the variety of the measurable evidence is directly visible
in the table. Deformation mechanisms are expressed as geometric
indicators in ‰ D, in mm. Leakage and durability mechanisms are
expressed as count-based indicators, e.g., defects per hundred square
meters. Source attribution is uneven across the two main codes. The
largest share of chains comes from CJJ/T 289, with the
DGTJ 08-2123 appendices and GB 50911 providing the rest. The single GB
50208 entry provides the leakage-grade definition.

\begin{longtable}[]{@{}
  >{\raggedright\arraybackslash}p{(\columnwidth - 8\tabcolsep) * \real{0.2074}}
  >{\raggedright\arraybackslash}p{(\columnwidth - 8\tabcolsep) * \real{0.1086}}
  >{\raggedright\arraybackslash}p{(\columnwidth - 8\tabcolsep) * \real{0.2373}}
  >{\raggedright\arraybackslash}p{(\columnwidth - 8\tabcolsep) * \real{0.2274}}
  >{\raggedright\arraybackslash}p{(\columnwidth - 8\tabcolsep) * \real{0.2193}}@{}}
\caption{Extracted FMEA knowledge base (rule layer), grouped by
degradation mechanism.}
\label{tab:12}\\
\toprule\noalign{}
\begin{minipage}[b]{\linewidth}\raggedright
\textbf{Mechanism category}
\end{minipage} & \begin{minipage}[b]{\linewidth}\raggedright
\textbf{Chains}
\end{minipage} & \begin{minipage}[b]{\linewidth}\raggedright
\textbf{Main component(s)}
\end{minipage} & \begin{minipage}[b]{\linewidth}\raggedright
\textbf{Indicator(s)}
\end{minipage} & \begin{minipage}[b]{\linewidth}\raggedright
\textbf{Source standard(s)}
\end{minipage} \\
\midrule\noalign{}
\endhead
\bottomrule\noalign{}
\endlastfoot
Water Leakage & 10 & Structural Connection, Waterproofing System &
L/(m\textsuperscript{2}$\cdot$d), L/d, count/100m\textsuperscript{2}, m\textsuperscript{2}, ratio (per 1000) & CJJ/T 289 Sec7; CJJ/T
289 Sec5 \\
Convergence Deformation & 8 & Segment Lining, Cut and Cover Structure &
L0 ratio, mm, ratio of {[}k{]}, ratio of {[}$\xi${]}, \textperthousand D & DGTJ 08-2123
Sec5; DGTJ 08-2123 App H \\
Longitudinal Deformation & 6 & Segment Lining, Cut and Cover Structure &
\% L, 1/m (curvature), mm, \textperthousand{} & GB 50911 Sec9-10-11; DGTJ 08-2123 App
H \\
Joint Dislocation & 5 & Structural Connection, Cut and Cover Structure &
mm & DGTJ 08-2123 App G; CJJ/T 289 Sec5 \\
Joint Opening & 5 & Structural Connection, Cut and Cover Structure & mm
& DGTJ 08-2123 App G; CJJ/T 289 Sec5 \\
Cracking & 4 & Segment Lining, Track Bed & mm & CJJ/T 289 Sec7; CJJ/T
289 Sec5 \\
Steel Bolt Corrosion & 3 & Segment Lining, Bolt & \% & CJJ/T 289 Sec5;
CJJ/T 289 Sec7 \\
Void Behind Lining & 3 & Track Bed, Segment Lining & m, m\textsuperscript{2}, mm & CJJ/T
289 Sec5; CJJ/T 289 Sec7 \\
Insufficient Capacity & 3 & Segment Lining, Bolt & m, ratio & DGTJ
08-2123 Sec5; CJJ/T 289 Sec5 \\
Material Degradation & 2 & Segment Lining & qualitative & CJJ/T 289
Sec5; CJJ/T 289 Sec7 \\
Spalling Peeling & 2 & Segment Lining & m\textsuperscript{2}, mm & CJJ/T 289 Sec5; DGTJ
08-2123 Sec5 \\
\textbf{Total} & \textbf{51} & 10 component categories & - & 7 standard
sections \\
\end{longtable}

Table \ref{tab:13} provides the decision for each of the twenty-five rings of the
two held-out tunnels, showing the convergence ratio and its graded
level, the joint dislocation and its graded level, the resulting ring
health level (and visualised in Figure \ref{fig:Fig10}), the assigned tier, and the
source clause by which each grade was assigned. Convergence, calculated
as the reduction of horizontal clearance and graded on the staggered
joint bands of CJJ/T 289, is between 1.59 and 4.22 ‰D across
Tunnel 8 and between 2.71 and 4.93 ‰D across Tunnel 9; the joint
dislocation, graded on Table 4, ranges from 1.6 to 5.4 mm and from 1.0
to 5.2 mm in the two tunnels. The ring health level is the more severe
of the two indicator grades, with each ring resolving to either level 1
or level 2 on the five-level scale for both tunnels.

In the provenance column, the local clause responsible for setting the
grade is recorded for each ring, so that each numeric decision in the
table is linked to the originating standard band and table, not to a
hard-coded threshold. In Tunnel 8, nine of the 13 rings are graded level
1 and four are level 2; in Tunnel 9, four of the 12 rings are level 1
and eight are level 2. Each ring in both tunnels is assigned a tier in
the acceptable-to-monitor range of the maintenance scale. When only the
deformation indicators are used, no ring is assigned to the
planned-repair or higher tiers.

\begin{longtable}[]{@{}
  >{\raggedright\arraybackslash}p{(\columnwidth - 16\tabcolsep) * \real{0.0803}}
  >{\raggedright\arraybackslash}p{(\columnwidth - 16\tabcolsep) * \real{0.0594}}
  >{\raggedright\arraybackslash}p{(\columnwidth - 16\tabcolsep) * \real{0.1098}}
  >{\raggedright\arraybackslash}p{(\columnwidth - 16\tabcolsep) * \real{0.0998}}
  >{\raggedright\arraybackslash}p{(\columnwidth - 16\tabcolsep) * \real{0.1294}}
  >{\raggedright\arraybackslash}p{(\columnwidth - 16\tabcolsep) * \real{0.1436}}
  >{\raggedright\arraybackslash}p{(\columnwidth - 16\tabcolsep) * \real{0.1056}}
  >{\raggedright\arraybackslash}p{(\columnwidth - 16\tabcolsep) * \real{0.1596}}
  >{\raggedright\arraybackslash}p{(\columnwidth - 16\tabcolsep) * \real{0.1125}}@{}}
\caption{Per-Ring decision with standard provenance.}
\label{tab:13}\\
\toprule\noalign{}
\begin{minipage}[b]{\linewidth}\raggedright
\textbf{Tunnel}
\end{minipage} & \begin{minipage}[b]{\linewidth}\raggedright
\textbf{Ring}
\end{minipage} & \begin{minipage}[b]{\linewidth}\raggedright
\textbf{Conv. (D)}
\end{minipage} & \begin{minipage}[b]{\linewidth}\raggedright
\textbf{Conv. lvl}
\end{minipage} & \begin{minipage}[b]{\linewidth}\raggedright
\textbf{Conv. band}
\end{minipage} & \begin{minipage}[b]{\linewidth}\raggedright
\textbf{Disloc. (mm)}
\end{minipage} & \begin{minipage}[b]{\linewidth}\raggedright
\textbf{Disloc. lvl}
\end{minipage} & \begin{minipage}[b]{\linewidth}\raggedright
\textbf{Ring health lvl}
\end{minipage} & \begin{minipage}[b]{\linewidth}\raggedright
\textbf{Tier}
\end{minipage} \\
\midrule\noalign{}
\endhead
\bottomrule\noalign{}
\endlastfoot
8 & 1 & 2.24 & 1 & $0 \le c \le 4$ & 3.0 & 1 & 1 & Acceptable \\
8 & 2 & 2.40 & 1 & $0 \le c \le 4$ & 5.4 & 2 & 2 & Monitor \\
8 & 3 & 2.17 & 1 & $0 \le c \le 4$ & 1.8 & 1 & 1 & Acceptable \\
8 & 4 & 1.59 & 1 & $0 \le c \le 4$ & 2.2 & 1 & 1 & Acceptable \\
8 & 5 & 3.38 & 1 & $0 \le c \le 4$ & 2.7 & 1 & 1 & Acceptable \\
8 & 6 & 2.94 & 1 & $0 \le c \le 4$ & 1.6 & 1 & 1 & Acceptable \\
8 & 7 & 2.48 & 1 & $0 \le c \le 4$ & 2.6 & 1 & 1 & Acceptable \\
8 & 8 & 2.29 & 1 & $0 \le c \le 4$ & 4.7 & 2 & 2 & Monitor \\
8 & 9 & 4.22 & 2 & $4 < c \le 6$ & 2.6 & 1 & 2 & Monitor \\
8 & 10 & 2.42 & 1 & $0 \le c \le 4$ & 3.1 & 1 & 1 & Acceptable \\
8 & 11 & 3.91 & 1 & $0 \le c \le 4$ & 3.3 & 1 & 1 & Acceptable \\
8 & 12 & 3.80 & 1 & $0 \le c \le 4$ & 4.3 & 2 & 2 & Monitor \\
8 & 13 & 3.36 & 1 & $0 \le c \le 4$ & 3.9 & 1 & 1 & Acceptable \\
9 & 1 & 4.60 & 2 & $4 < c \le 6$ & 2.5 & 1 & 2 & Monitor \\
9 & 2 & 3.73 & 1 & $0 \le c \le 4$ & 3.9 & 1 & 1 & Acceptable \\
9 & 3 & 4.93 & 2 & $4 < c \le 6$ & 2.5 & 1 & 2 & Monitor \\
9 & 4 & 4.07 & 2 & $4 < c \le 6$ & 3.4 & 1 & 2 & Monitor \\
9 & 5 & 3.57 & 1 & $0 \le c \le 4$ & 5.2 & 2 & 2 & Monitor \\
9 & 6 & 4.50 & 2 & $4 < c \le 6$ & 3.9 & 1 & 2 & Monitor \\
9 & 7 & 3.88 & 1 & $0 \le c \le 4$ & 1.0 & 1 & 1 & Acceptable \\
9 & 8 & 3.20 & 1 & $0 \le c \le 4$ & 3.5 & 1 & 1 & Acceptable \\
9 & 9 & 2.71 & 1 & $0 \le c \le 4$ & 4.0 & 2 & 2 & Monitor \\
9 & 10 & 3.66 & 1 & $0 \le c \le 4$ & 2.5 & 1 & 1 & Acceptable \\
9 & 11 & 3.60 & 1 & $0 \le c \le 4$ & 4.2 & 2 & 2 & Monitor \\
9 & 12 & 4.34 & 2 & $4 < c \le 6$ & 3.0 & 1 & 2 & Monitor \\
\end{longtable}


\begin{figure}[htbp]
  \centering
  \includegraphics[width=0.92\textwidth]{Fig_PDF/Picture10.pdf}
  \caption{ Distribution of ring health levels along the tunnels}
  \label{fig:Fig10}
\end{figure}


The ring-level results of the tunnel scale are summarised in Table \ref{tab:14}
and presented against the vision channel leakage evidence. In both
tunnels, the worst deformation level is 2, which is consistent with the
ring distributions in Table B. The leakage channel differs between the
two tunnels. Tunnel 8 is graded at leakage level 4 and has a
linear-leakage class of an estimated 4.0 m 2. Tunnel 9 is graded at
leakage level 2 and has a seepage class of an estimated 0.94
m\textsuperscript{2}.

The overall tire for Tunnel 8 is given as level 4, Action Required by
its leakage grade, following the max-severity rule, which combines the
structural and leakage channels. The overall tier for Tunnel 9 is at
level 2, Monitor, which is consistent with its worst deformation level
and its leakage level. The governing-channel column directly records
this difference, associating the outcome of the Tunnel 8 with the
leakage channel and the outcome of the Tunnel 9 with a deformation grade
that is not exceeded by the leakage channel.

\begin{longtable}[]{@{}
  >{\raggedright\arraybackslash}p{(\columnwidth - 14\tabcolsep) * \real{0.0803}}
  >{\raggedright\arraybackslash}p{(\columnwidth - 14\tabcolsep) * \real{0.0798}}
  >{\raggedright\arraybackslash}p{(\columnwidth - 14\tabcolsep) * \real{0.1695}}
  >{\raggedright\arraybackslash}p{(\columnwidth - 14\tabcolsep) * \real{0.1496}}
  >{\raggedright\arraybackslash}p{(\columnwidth - 14\tabcolsep) * \real{0.1595}}
  >{\raggedright\arraybackslash}p{(\columnwidth - 14\tabcolsep) * \real{0.0997}}
  >{\raggedright\arraybackslash}p{(\columnwidth - 14\tabcolsep) * \real{0.1296}}
  >{\raggedright\arraybackslash}p{(\columnwidth - 14\tabcolsep) * \real{0.1320}}@{}}
\caption{Tunnel-Level decision summary.}
\label{tab:14}\\
\toprule\noalign{}
\begin{minipage}[b]{\linewidth}\raggedright
\textbf{Tunnel}
\end{minipage} & \begin{minipage}[b]{\linewidth}\raggedright
\textbf{Rings}
\end{minipage} & \begin{minipage}[b]{\linewidth}\raggedright
\textbf{Ring health levels (count)}
\end{minipage} & \begin{minipage}[b]{\linewidth}\raggedright
\textbf{Worst deformation level}
\end{minipage} & \begin{minipage}[b]{\linewidth}\raggedright
\textbf{Leakage class (area)}
\end{minipage} & \begin{minipage}[b]{\linewidth}\raggedright
\textbf{Leakage level}
\end{minipage} & \begin{minipage}[b]{\linewidth}\raggedright
\textbf{Overall tier}
\end{minipage} & \begin{minipage}[b]{\linewidth}\raggedright
\textbf{Governing channel}
\end{minipage} \\
\midrule\noalign{}
\endhead
\bottomrule\noalign{}
\endlastfoot
8 & 13 & L1×9, L2×4 & 2 & Linear Leakage (4.0 m²) & 4 & \textbf{Action
Required (L4)} & Leakage \\
9 & 12 & L1×4, L2×8 & 2 & Seepage (0.94 m²) & 2 & \textbf{Monitor (L2)}
& Deformation \\
\end{longtable}

\section{Conclusion}\label{conclusion}

The goal of this study was to close the loop that most tunnel
digital-twin frameworks leave open, carrying raw point-cloud and image
evidence through to a standard referenced maintenance prescription
rather than stopping at detection or forecast. The pipeline integrates a
self-supervised point-cloud segmentation stage, a geometric
deformation-reconstruction stage, and an ontology-based decision layer.
The pipeline was exercised end-to-end on two fully held-out
shield-tunnel sections of the Nanjing Metro, with an Australian
reference structure provided as supplementary training data.

The best performing configuration for leakage perception was full
fine-tuning of the Sonata encoder on the in-domain S3DIS leakage data at
full supervision (10\% of the total scanned dataset), with a leakage IoU
of 94.34, mean IoU of 93.90, and mean accuracy of 98.60, with an overall
accuracy and macro-F1 of 0.97 and a Cohen's $\kappa$ and MCC of 0.95. The given
fine-tuned model performed better than the strongest external
point-cloud baseline on the same, but the whole of the benchmark, the
BBE-Net ensemble, on both leakage IoU and mean accuracy. At the class
level, the leakage region recovered at an IoU of 0.945 and a recall of
0.943. Segments, pockets, and joints were segmented at IoUs of 1.000,
0.974, and 0.840, respectively, the lower joint value reflecting the
small fraction of true-leakage points that the confusion matrix assigned
to the adjacent joint class.

The experiments on data combination and data fraction provided useful
results for practical annotation. In all matched settings, full
fine-tuning beat linear probing. The in-domain S3DIS data with abundant
labels resulted in the best performance. The encoder was free to adapt,
with the S3DIS-only model at full supervision (94.34 leakage IoU)
standing on top of its S3DIS -ART counterpart (92.17). The additional
Australian data demonstrated its performance in the setting with scarce
labels and a frozen encoder, improving the 25\% linear-probe leakage IoU
from 45.41 to 70.02 and the corresponding accuracy from 0.69 to 0.83.
The picture for label-efficiency is correspondingly favourable: full
fine-tuning retained an accuracy of 0.93 and a leakage IoU of 87.59 with
only a quarter of the labels, and since the annotated pool was itself
one-tenth of the scanned corpus, that quarter-budget corresponds to
about 2.5\% of the available scan data. Thus, strong leakage
segmentation can be obtained at a small fraction of the labelling effort
assumed by previous work, with cross-site auxiliary data acting as a
surrogate for labels precisely in the low-annotation regime where it is
most needed.

The deformation stage reduced each reconstructed ring to a stable set of
indicators across both held-out tunnels. Clearance convergence ranged
from 1.59 to 4.22 ‰D in Tunnel 8 and from 2.71 to 4.93 ‰D in Tunnel 9,
with ovalization between 40 to 52 mm, maximum joint dislocation at 5.4
mm, and joint rotation from 2° to 9°, except for two rings that reached
about 12°. The ovalization, the most sensitive indicator, was calculated
through multi-zone polynomial reconstruction and a raw conic fit,
resulting in a mean absolute difference of 1.37 mm (1.15 mm in Tunnel 8
and 1.62 mm in Tunnel 9), with a maximum discrepancy of 3.6 mm for the
highest-ovalization ring. This method reduced the impact of acquisition
errors, ensuring that systematic offsets in input geometry do not skew
the structural grade, thus validating the deformation evidence through
cross-method agreement and consistent measurements.

The decision layer converted this perceptual and geometric evidence into
prescriptions that were traceable and referenced to standards.
Schema-guided extraction populated 51 FMEA reasoning chains from seven
sections of the governing Chinese urban-rail codes over eleven
degradation mechanisms. Each of the twenty-five rings was graded by
traversing these chains against its measured indicators; the fired band
and originating clause were retained for each decision. From the
deformation indicators alone, all rings were resolved to health levels 1
and 2, the acceptable-to-monitor range. The leakage channel then
separated the two tunnels, promoting Tunnel 8 to an Action-Required tier
on a linear-leakage class over an estimated 4.0 m², while Tunnel 9
remained at Monitor under a smaller seepage area. Each grade is taken
from the populated rule base rather than a hard-coded threshold, so a
superseding standard or different tunnel type can be accommodated by
revising the rule base alone, and every prescription remains inspectable
back to the clause that governs it.

Many additional mechanisms manifest visually rather than geometrically,
making them difficult to capture from point-cloud data alone. Extending
the methodology to image-based foundation models for photographic
evidence of cracks and surface defects, in a self-supervised and
label-efficient manner, would complement the geometric analysis and
align with the multimodal design intent.

The knowledge layer is currently based on Chinese urban-rail standards,
using minimal information for each stage. Integrating maintenance and
acceptance guidelines from other jurisdictions would enhance its
applicability and test the extraction against various regulatory
vocabularies without changing the underlying pipeline.

Lastly, the assessment relies on a single scan epoch, reflecting
deviation from nominal design geometry rather than in-service changes.
Incorporating periodic scanning campaigns would transform these static
grades into deformation trajectories, enabling tracking of changes over
time and moving toward predictive maintenance scheduling.

\section*{Acknowledgements}\label{acknowledgements}
\addcontentsline{toc}{section}{Acknowledgements}

The authors acknowledge the support of the asset client for facilitating
access to the case study tunnel for the multimodal sensing campaign, and
of industry partners who contributed to validation of the FMEA
extraction results.

\section*{Declaration}\label{declaration}
\addcontentsline{toc}{section}{Declaration}

The authors declare no competing interests. Large language models (LLM)
were used in this work in two distinct capacities. As part of the
research methodology, Anthropic Claude (Claude Opus 4.8) was used for
schema-guided extraction of FMEA reasoning chains from maintenance
standards under expert validation; this use is described in full in
Section X, including model version. Separately, a large language model
was used to assist with language editing and readability during
manuscript preparation. After using these tools, the authors reviewed
and edited all such content as needed and took full responsibility for
the content of the publication.

\section*{Data availability}\label{data-availability}
\addcontentsline{toc}{section}{Data availability}

The scanning dataset from Australian tunnel is subject to operator
confidentiality and cannot be released. The S3DIS dataset is available
online from:

\url{https://www.kaggle.com/datasets/yueqianshen/s3dis-leakage}


\clearpage
\bibliographystyle{elsarticle-harv}
\nocite{*}
\bibliography{references_elsevier}

\clearpage
\section*{Appendix A- Per-Ring deformation
cross-section}\label{appendix-a--per-ring-deformation-cross-section}
\addcontentsline{toc}{section}{Appendix A- Per-Ring deformation
cross-section}


\begin{figure}[H]
  \centering
  \includegraphics[width=0.92\textwidth]{Fig_PDF/PictureA1.pdf}
  \caption{ Per-ring radial deformation for tunnel 8}
  \label{fig:FigA1}
\end{figure}


\begin{figure}[H]
  \centering
  \includegraphics[width=0.92\textwidth]{Fig_PDF/PictureA2.pdf}
  \caption{ Per-ring radial deformation for tunnel 9}
  \label{fig:FigA2}
\end{figure}

\section*{Appendix C- FMEA-chain extraction
pipeline}\label{appendix-c--fmea-chain-extraction-pipeline}
\addcontentsline{toc}{section}{Appendix C- FMEA-chain extraction
pipeline}

The pseudocode below describes the schema-guided extraction pipeline
used to build the FMEA knowledge base of Section {[}X{]}. The standards
corpus is read and split into fixed-size character chunks, each chunk is
sent to a large language model under a system prompt that injects the
controlled vocabulary and a strict JSON schema, the returned JSON array
is parsed, and all parsed records are accumulated into the knowledge
base together with a per-chunk timing and provenance log. Generation is
deterministic (temperature 0, fixed seed).

\begin{algorithm}[H]
\caption{Schema-guided FMEA-chain extraction with optional vision perception.}
\label{alg:fmea}
\begin{algorithmic}[1]
\Require Ordered standards corpus
  $\mathcal{S}=[(\mathit{file}_1,\mathit{label}_1),\dots]$;
  controlled vocabulary $V$ (component, mechanism, health-level, indicator,
  intervention tokens); output schema $\Sigma$; language model $\mathcal{M}$
  (Claude Opus~4.8); chunk size $c$ (characters); decoding parameters
  $\theta=(\text{temperature}=0,\ \text{seed},\ \text{context window})$.
\Ensure FMEA-chain set $\mathcal{F}$ and run log $\mathcal{L}$.
\State $\mathcal{F}\gets[\,]$;\quad $\mathcal{L}\gets[\,]$;\quad
  $\mathit{Chunks}\gets[\,]$
\State $P_{\text{sys}}\gets \Call{BuildSystemPrompt}{\Sigma,V}$
  \Comment{controlled vocabulary injected verbatim}
\For{each $(\mathit{file},\mathit{label})\in\mathcal{S}$}
  \Comment{Step 1: read and chunk}
  \State $\mathit{text}\gets\Call{ReadUTF8}{\mathit{file}}$;\quad
    $\mathit{start}\gets 0$;\quad $k\gets 1$
  \While{$\mathit{start}<\Call{Length}{\mathit{text}}$}
    \State $\mathit{end}\gets \mathit{start}+c$
    \If{$\mathit{end}<\Call{Length}{\mathit{text}}$}
      \State $\mathit{end}\gets\Call{NextNewlineAtOrAfter}{\mathit{text},\mathit{end}}$
        \Comment{break on a line boundary}
    \EndIf
    \State $\mathit{chunk}\gets\Call{Trim}{\mathit{text}[\mathit{start}:\mathit{end}]}$
    \If{$\mathit{chunk}$ is non-empty}
      \State append $(\mathit{label},k,\mathit{chunk})$ to $\mathit{Chunks}$;\quad
        $k\gets k+1$
    \EndIf
    \State $\mathit{start}\gets \mathit{end}+1$
  \EndWhile
\EndFor
\For{each $(\mathit{label},k,\mathit{chunk})\in\mathit{Chunks}$}
  \Comment{Step 2: extract per chunk}
  \State $P_{\text{user}}\gets\Call{FormatUserPrompt}{\mathit{label},k,\mathit{chunk}}$
  \State $\mathit{raw}\gets\mathcal{M}\bigl(P_{\text{sys}},P_{\text{user}};
    \theta,\text{format}=\text{JSON}\bigr)$
    \Comment{schema-constrained generation}
  \State $R\gets\Call{ParseJSONArray}{\mathit{raw}}$
    \Comment{strip fences; tolerant parse}
  \If{$R$ is not a list}
    \State $R\gets[R]$ if $R$ is a non-empty object, else $R\gets[\,]$
  \EndIf
  \State append all records in $R$ to $\mathcal{F}$
  \State append $(\mathit{label},k,|R|,\mathit{elapsed})$ to $\mathcal{L}$
\EndFor
\If{vision is enabled}
  \Comment{Step 3: optional RGB perception}
  \For{each image $I$ in $\mathit{ImageDir}$}
    \State $r\gets\mathcal{M}_{\text{vision}}\bigl(P_{\text{sys}}^{\text{vision}},
      \Call{FormatVisionPrompt}{I};\theta\bigr)$
    \If{$r$ is an object}
      \State append $r$ to $\mathcal{F}_{\text{image}}$
    \EndIf
    \State append $(\text{image name},\mathit{elapsed},\text{status})$ to $\mathcal{L}$
  \EndFor
\EndIf
\State \Call{Save}{$\mathcal{F},\mathcal{F}_{\text{image}},\mathcal{L}$}
  \Comment{JSON knowledge base + run log}
\State \Return $(\mathcal{F},\mathcal{L})$
\end{algorithmic}
\end{algorithm}

\textbf{Notes on the schema constraint.} The system prompt fixes the
domain to segmental shield metro linings and the authoritative severity
scale to the five-level health degree of CJJ/T 289 (Level 1 best to
Level 5 worst). The model is instructed to (i) return only a JSON array
with no fields outside the schema, (ii) use only controlled-vocabulary
tokens and otherwise emit null rather than invent a token, (iii) copy
graded numeric bands verbatim into grade\_bands with their unit and
source clause, never paraphrasing a number, (iv) map qualitative leakage
descriptions to the leakage-class tokens and health levels while
recording the source clause, and (v) return an empty array for chunks
that contain no FMEA-relevant content. Each emitted record carries its
source\_documents label and per-threshold source\_reference, so that
every chain in the resulting knowledge base is traceable to the clause
or table it was drawn from.


\end{document}
